\documentclass[11pt]{article}
\usepackage[margin=1in]{geometry}
\usepackage{amsmath,amssymb,amsthm,mathtools}
\usepackage[hidelinks,pdfusetitle]{hyperref}
\usepackage{microtype}
\usepackage{array}
\usepackage{etoolbox}
\newcolumntype{L}[1]{>{\raggedright\arraybackslash}p{#1}}
\newcommand{\DOI}[1]{\href{https://doi.org/#1}{\nolinkurl{#1}}}
% Identifiers of the strong-coupling companion preprint and of the present note.
% The release step fills \SELFDOI.
\newcommand{\STRONGDOI}{10.5281/zenodo.24414616}
\newcommand{\SELFDOI}{}
\newcommand{\record}[2]{\ifdefempty{#1}{#2}{\DOI{#1}}}
\theoremstyle{plain}
\newtheorem{theorem}{Theorem}
\newtheorem{corollary}{Corollary}[theorem]
\newtheorem{proposition}[theorem]{Proposition}
\newtheorem{lemma}{Lemma}
\newtheorem{remark}{Remark}
\theoremstyle{definition}
\newtheorem{definition}{Definition}
\DeclareMathOperator{\Tr}{Tr}
\newcommand{\Nodd}{N_{\mathrm{odd}}}
\newcommand{\Dop}{\Delta}
\newcommand{\ad}{\operatorname{ad}}
\setlength{\emergencystretch}{2em}
\title{Complex-rotation pairing bounds for a cold-projected bipolaron-phonon
effective model and obstructions to microscopic extensions in the
two-dimensional Holstein model}
\author{Davis Rens\\SuperMatics Inc.\\\href{https://orcid.org/0000-0001-5839-5197}{ORCID 0000-0001-5839-5197}}
\date{}
\begin{document}
\maketitle

\begin{abstract}
Complex rotations \cite{bfu} are the standard device that converts a uniform
rotation of the charge into algebraic rather than exponential decay of
correlations in two dimensions. Applied to the spinful
Holstein model at large $U/\Omega$, they meet a specific difficulty: the pairing
field is charge two, so a rotation that assigns a slowly varying phase to the
electronic charge acts on a virtual pair as soon as the pair is created, and the
one-hop mixing that a Lang--Firsov transformation leaves behind is what the
rotation must undo. This note works out what that difficulty costs. We take the
Lang--Firsov plus first-order Schrieffer--Wolff transform of the Holstein
Hamiltonian on the square torus, project it onto the cold manifold
$\Nodd=0$, $n^{\mathrm{ph}}_z\le M=\lfloor U/(64\Omega)\rfloor$, and apply a
$U(1)$ charge rotation to the projected effective Hamiltonian. For
$x=t/U\le2^{-14}$ and $X=U/\Omega\ge1$ we prove that the polaron-pair correlation
obeys $\bigl|\langle\Dop^\dagger_{x_0}\Dop_{y_0}\rangle_{\mathrm{cold}}\bigr|
\le(d+1)^{-\xi_{\mathrm{cold}}(\beta)}$ with an explicit exponent
$\xi_{\mathrm{cold}}(\beta)=4-16A(\beta)$ for $A(\beta)<1/8$ and
$\xi_{\mathrm{cold}}(\beta)=1/(4A(\beta))$ otherwise, where
$A(\beta)=8\beta UK_4(x)e^{-7U/(8\Omega)}$ and $K_4(x)\le2^{20}x^2$. Hence the
correlations are summable in two dimensions whenever
$\tau>2^{26}\,(t^2/U)\,e^{-7U/(8\Omega)}$, and the same exponent governs the
original equal-time electronic pair field with prefactor $4$.

The narrowed ceiling sits below the released $2^{27}t^{2}/U$ bound
\cite{rens2026a} for every $U/\Omega\ge1$, but it is a statement about an
effective model and not about the full Holstein Gibbs state. We therefore also
prove two obstructions that locate the boundary of the present method. First, a
rotation applied directly to the microscopic Hamiltonian meets a floor at the
scale $t^2/U$: on two sites with $N=2$, $S_z=0$, virtual neutral excursions lower the
Hermitian-twist free energy by at least
$\tfrac{1}{400}(t^2/U)\sinh^2(\delta/2)$ for all $\Omega>0$, $\beta U\ge4$ and
$0<t/U\le1/16$, and the zero-temperature Franck--Condon resolvent obeys
$\sum_n e^{-X}X^n/(n!\,(U+n\Omega))\ge1/(2U)$ for all $\Omega>0$. Second, the
Gibbs weight of the cold manifold vanishes like
$\exp(-|\Lambda|e^{-\beta\Omega(M+1)})$ at fixed $\tau/\Omega$, so the projection
itself gives nothing at fixed temperature once the phonon occupations above $M$
are admitted. Control of the full Gibbs state in that regime remains open.

The bounds are on a clustering temperature of one two-point function inside an
effective state. They are not bounds on a critical temperature, no phase
transition is proved or excluded, and nothing here asserts that a material is a
superconductor. A finite certificate program recomputes every constant; it
produces a certified replay, rejects sixteen corrupted variants of the claimed
inequalities one at a time, and reproduces the two-site numerical checks.
\end{abstract}

\section{What this note does, and what it does not}

For the two-dimensional Holstein model two upper bounds on the clustering
temperature of the electronic pair field are released: the bound
$k_{\mathrm{B}}T_{\mathrm{pair}}\le2^{27}t^{2}/U$ \cite{rens2026a} and the
sharper strong-coupling bound $k_{\mathrm{B}}T_{\mathrm{pair}}\le6000\,t^{2}/U$
(improving to $4500\,t^{2}/U$ on tori of even side) of the companion preprint
\cite{rens2026strong}. Both hold for the full Gibbs state, impose no restriction
on the ratio $t/U$, are uniform in the phonon frequency $\Omega>0$, in the torus
side $L\ge4$ and in the chemical potential $\mu\in\mathbb{R}$, and prove
exponential decay by expansions in electronic trajectories rather than in the
electronic charge. The question addressed here is what a complex rotation of the
electronic charge, the classical device for algebraic rather than exponential
decay in two dimensions \cite{bfu}, does in the same model at large $U/\Omega$.

The answer has three parts, and the negative parts are as much of a result as
the positive one. The rotation acts on the pairing field through the charge
carried by the intermediate states of the electron-phonon cloud. In the
Lang--Firsov frame the charge-carrying terms of the model are single hops dressed
by displacement operators, of norm $t$; a rotation that multiplies each hop by a
phase is therefore reimbursed, in the twisted trace, by a factor that grows with
the twist. That growth is what a cluster expansion removes by summing over
connected sets of hops, and on the cold manifold, where every hop changes the
number of singly occupied sites, the summation produces two independent
displacement operators with displacement magnitude at least $\sqrt{2X}$. Those
two operators are what buy the factor $e^{-7U/(8\Omega)}$. Off the cold
manifold, a single virtual pair excursion across one bond creates and destroys
the pair on the same bond, the two displacement operators cancel, and the
suppression disappears. This is the content of Proposition~\ref{prop:floor}: on
two sites the free-energy gain of the twisted Hamiltonian is bounded below by a
$t^{2}/U$ term that no choice of rotation removes.

The note is organised so that each statement can be read on its own. Section~2
sets up the model, the cold manifold and the statements.
Section~3 proves the cold-projected theorem and its corollary.
Section~4 proves the microscopic obstruction and records the numerical
cross-check. Section~5 proves the thermal statements about the twisted pair
hopping and about the Gibbs weight of the cold manifold. Section~6 states the
limitations. Section~7 describes the certificate program.

Throughout, a statement is called proved when a written derivation is given
here, certified when a finite arithmetic identity or inequality on which it rests
is recomputed in exact rational or outward-rounded interval arithmetic by the
certificate program, and numerical when it is checked by diagonalisation of a
finite matrix. Two items are called open, and they are named as open where they
occur.

\subsection{What is not claimed}

The model of Theorem~\ref{thm:main} is not the Holstein model. It is the
cold-projected effective model obtained by conjugating the Holstein Hamiltonian
with a Lang--Firsov transformation, truncating the resulting Schrieffer--Wolff
series at first order, and projecting onto the sector of doubly or zero occupied
sites with phonon occupations bounded by $M=\lfloor U/(64\Omega)\rfloor$.
Section~2.3 states the construction and the projection exactly. The projection is
a genuine restriction: the state in which the correlations of
Theorem~\ref{thm:main} are computed is the Gibbs state of the projected
Hamiltonian, not the Gibbs state of the Holstein Hamiltonian, and
Proposition~\ref{prop:thermal}(2) shows that the projection cannot be undone by
restricting the trace.

The bounds are on the clustering temperature of one two-point function, in the
sense of Definition~\ref{def:tpair}. They are not bounds on a critical
temperature. No phase transition is proved or excluded at any temperature, and
nothing here says that the model, at any temperature, is a superconductor or
shows any form of long-range order. The finite certificate program proves
arithmetic, not analysis; it does not replace any operator inequality, each of
which has a written derivation in Sections~3 to 5.

\section{Setting and statements}\label{sec:model}

\subsection{The spinful Holstein model on a square torus}\label{sec:conv}

Fix an integer side $L\ge4$ and let $\Lambda=T_L=(\mathbb{Z}/L\mathbb{Z})^{2}$
be the square torus with its nearest-neighbour edge set
$E(\Lambda)=\{\{x,y\}\subset\Lambda:|x-y|_{1}=1\}$, where
$|x-y|_{1}$ is the torus distance in the $\ell_{1}$ metric, written $d(x,y)$
below. Each site carries one oscillator coordinate and a spinful electronic
level with four states. With $t\ge0$, $\Omega>0$ (the phonon frequency, in units
where $\hbar=1$), $\alpha=\sqrt{U/(2\Omega)}$, $g=\sqrt{U\Omega/2}$ (the
electron-phonon coupling in the convention $g(a^\dagger_x+a_x)(n_x-1)$) and real
$\mu$, the Hamiltonian of \cite{holstein1959} is
\begin{equation}\label{eq:H}
H_\Lambda=-t\sum_{e\in E(\Lambda)}\sum_{\sigma}B_{e,\sigma}
+\Omega\sum_{x\in\Lambda}a_x^\dagger a_x
-\sqrt{\frac{U\Omega}{2}}\sum_{x\in\Lambda}(n_x-1)(a_x^\dagger+a_x)
-\mu\sum_{x\in\Lambda}n_x ,
\end{equation}
with $B_{e,\sigma}=c^\dagger_{u\sigma}c_{v\sigma}+c^\dagger_{v\sigma}c_{u\sigma}$
for $e=\{u,v\}$ and $\sigma\in\{\uparrow,\downarrow\}$,
$n_x=n_{x\uparrow}+n_{x\downarrow}$, and $U>0$ the polaronic energy
$U=F^{2}/K$ of the oscillator. There is no direct electron-electron interaction.
We write $\tau=k_{\mathrm{B}}T=\beta^{-1}>0$, and the two dimensionless ratios
\begin{equation}\label{eq:ratios}
x=\frac{t}{U}\le2^{-14},\qquad X=\frac{U}{\Omega}\ge1,
\end{equation}
are the only small or large parameters used in Theorem~\ref{thm:main}.
The projector onto singly occupied sites is
$\mathbf{1}_{\{n_x=1\}}=1-(n_x-1)^{2}$ on the four electronic states, and
\begin{equation}\label{eq:Nodd}
\Nodd=\sum_{x\in\Lambda}\bigl(1-(n_x-1)^{2}\bigr),\qquad
N=\sum_{x\in\Lambda}n_x .
\end{equation}
The equal-time on-site pair field is $\Dop_x=c_{x\downarrow}c_{x\uparrow}$, with
the order of the two operators fixed once and for all.

\begin{definition}[Clustering temperatures]\label{def:tpair}
Let $T_{\mathrm{pair}}$ be the infimum of the temperatures $\tau$ with the
following property: for every $\tau'>\tau$, every $L\ge4$, every real $\mu$ and
every $\Omega>0$ there are constants $A_0>0$ and $\xi>0$, independent of $L$,
$\mu$ and $\Omega$, such that
$\bigl|\langle\Dop_x^\dagger\Dop_y\rangle\bigr|\le A_0e^{-d(x,y)/\xi}$ for all
distinct $x,y\in\Lambda$ and
$\sup_{L\ge4}\max_x\sum_{y}|\langle\Dop_x^\dagger\Dop_y\rangle|<\infty$, the
expectation being that of $H_\Lambda$. For fixed $U>0$, $t\ge0$ and $\Omega>0$
with $X=U/\Omega\ge1$, let $T_{\mathrm{pair}}^{\mathrm{cold}}$ be the infimum of
the temperatures $\tau$ such that for every $\tau'>\tau$ there are constants
$A_0>0$ and $\xi>2$, independent of $L\ge4$ and $\mu\in\mathbb{R}$, with
$\bigl|\langle\Dop_x^\dagger\Dop_y\rangle_{\mathrm{cold}}\bigr|\le A_0(d(x,y)+1)^{-\xi}$
for all distinct $x,y\in\Lambda$ and
$\sup_{L\ge4}\max_x\sum_{y}|\langle\Dop_x^\dagger\Dop_y\rangle_{\mathrm{cold}}|<\infty$,
where $\langle\cdot\rangle_{\mathrm{cold}}$ is the state \eqref{eq:coldstate}.
\end{definition}

$T_{\mathrm{pair}}$ and $T_{\mathrm{pair}}^{\mathrm{cold}}$ are clustering
temperatures: upper thresholds for uniform row summability of one two-point
function (with exponential decay uniform in $\Omega$ for $T_{\mathrm{pair}}$, and
algebraic decay with exponent above $2$ at fixed $X=U/\Omega\ge1$ for
$T_{\mathrm{pair}}^{\mathrm{cold}}$). They say nothing about a phase transition,
and they are not critical temperatures. The definition of $T_{\mathrm{pair}}$ is
the one used in \cite{rens2026a,rens2026strong}; $T_{\mathrm{pair}}^{\mathrm{cold}}$
is the corresponding summable-clustering threshold for the effective model of
Section~2.3, with its dependence on $M$ carried by $X$ through
$M=\lfloor X/64\rfloor$.

\subsection{Lang--Firsov frame and first-order Schrieffer--Wolff generator}\label{sec:lf}

Let $D_x(\gamma)=\exp\bigl(\gamma(a_x^\dagger-a_x)\bigr)$ for real $\gamma$ be the
Weyl displacement, so that $D_x(-\gamma)a_xD_x(\gamma)=a_x+\gamma$, and let the
Lang--Firsov transformation \cite{hohenadler2007} be
\begin{equation}\label{eq:ulf}
U_{\mathrm{LF}}=\exp\Bigl(-\alpha\sum_{x\in\Lambda}(n_x-1)(a_x^\dagger-a_x)\Bigr),
\qquad \alpha=\sqrt{\frac{X}{2}} .
\end{equation}
Since $U_{\mathrm{LF}}$ commutes with $n_x$ and therefore with the charge,
\begin{equation}\label{eq:lfacts}
U_{\mathrm{LF}}a_xU_{\mathrm{LF}}^\dagger=a_x+\alpha(n_x-1),\qquad
U_{\mathrm{LF}}c^\dagger_{x\sigma}U_{\mathrm{LF}}^\dagger=c^\dagger_{x\sigma}D_x(-\alpha),\qquad
U_{\mathrm{LF}}c_{x\sigma}U_{\mathrm{LF}}^\dagger=c_{x\sigma}D_x(+\alpha).
\end{equation}
With $\Omega\alpha^{2}=U/2$ and $(n_x-1)^{2}=1-n_{x,\mathrm{odd}}$ the linear
coupling cancels:
$\Omega a^\dagger_xa_x-\tfrac{U}{2}(n_x-1)^{2}=\Omega n^{\mathrm{ph}}_x+\tfrac{U}{2}n_{x,\mathrm{odd}}-\tfrac{U}{2}$.
Dropping the additive constant $-\tfrac{U}{2}|\Lambda|$,
$H_{\mathrm{LF}}=U_{\mathrm{LF}}H_\Lambda U_{\mathrm{LF}}^\dagger+\tfrac{U}{2}|\Lambda|$ reads
\begin{equation}\label{eq:HLF}
H_{\mathrm{LF}}=D_0+H_{\mathrm{ph}}+V_{\mathrm{mix}}+V_{\mathrm{diag}},\qquad
D_0=\frac{U}{2}\Nodd-\mu N,\qquad H_{\mathrm{ph}}=\Omega\sum_{x\in\Lambda}n^{\mathrm{ph}}_x ,
\end{equation}
where, fixing a reference orientation $u\leftarrow v$ of each undirected edge
$e=\{u,v\}$ and writing $A_e=\sum_\sigma c^\dagger_{u\sigma}c_{v\sigma}$,
$P_{e,0}$ and $P_{e,2}$ for the projectors onto $n_{u,\mathrm{odd}}+n_{v,\mathrm{odd}}=0$
and $2$, $A_{e,\mathrm{mix}}=P_{e,0}A_eP_{e,2}+P_{e,2}A_eP_{e,0}$ and
$A_{e,\mathrm{diag}}=A_e-A_{e,\mathrm{mix}}$,
\begin{equation}\label{eq:Vmix}
V_{\mathrm{mix}}=\sum_{e\in E(\Lambda)}V_{e,\mathrm{mix}},\quad
V_{e,\mathrm{mix}}=-t\bigl(A_{e,\mathrm{mix}}D_u(-\alpha)D_v(\alpha)+\mathrm{h.c.}\bigr),
\end{equation}
and $V_{e,\mathrm{diag}}$ is given by the same expression with
$A_{e,\mathrm{diag}}$ in place of $A_{e,\mathrm{mix}}$, with
$V_{\mathrm{diag}}=\sum_eV_{e,\mathrm{diag}}$. On the sixteen electronic states of
the two sites of a bond one has
\begin{equation}\label{eq:Anorms}
\|A_{e,\mathrm{mix}}\|_{\mathrm{op}}=\sqrt{2}, \qquad
\|A_{e,\mathrm{diag}}\|_{\mathrm{op}}=1 .
\end{equation}
The skew-Hermitian generator of the first-order Schrieffer--Wolff
transformation is
\begin{equation}\label{eq:S1}
S_1=\sum_{e\in E(\Lambda)}S_{1,e},\qquad
S_{1,e}=x\Bigl[\bigl(P_{e,0}A_eP_{e,2}-P_{e,2}A_eP_{e,0}\bigr)D_u(-\alpha)D_v(\alpha)-\mathrm{h.c.}\Bigr].
\end{equation}
Three elementary facts are used below. First,
$\|S_{1,e}\|_{\mathrm{op}}\le2\sqrt{2}\,x$ on the full Hilbert space, because
$D_x(\pm\alpha)$ is unitary. Second, $S_1$ commutes with $N$, and
$\bigl[\tfrac{U}{2}\Nodd,P_{e,0}A_eP_{e,2}\bigr]=-UP_{e,0}A_eP_{e,2}$, so
\begin{equation}\label{eq:cancel}
[S_1,D_0]=-V_{\mathrm{mix}}
\end{equation}
identically on $\mathcal{H}_\Lambda$. Third, both $S_1$ and $H_{\mathrm{ph}}$
commute with $N$, so the conjugate of $H_\Lambda$ by the charge group below is
again $H_{\mathrm{LF}}$ with rescaled hop amplitudes.

\subsection{The cold manifold and its effective Hamiltonian}\label{sec:cold}

For $X\ge1$ put
\begin{equation}\label{eq:M}
M=\Bigl\lfloor\frac{X}{64}\Bigr\rfloor\in\mathbb{Z}_{\ge0}
\end{equation}
and define the orthogonal projector
\begin{equation}\label{eq:pcold}
P_{\mathrm{cold}}=\prod_{z\in\Lambda}\Bigl((1-n_{z,\mathrm{odd}})\,\mathbf{1}_{\{n^{\mathrm{ph}}_z\le M\}}\Bigr),
\qquad \mathcal{H}_{\mathrm{cold}}=\operatorname{ran}P_{\mathrm{cold}} .
\end{equation}
Then $\dim\mathcal{H}_{\mathrm{cold}}=[2(M+1)]^{|\Lambda|}<\infty$,
$P_{\mathrm{cold}}$ commutes with every site charge $S_z=\tfrac{1}{2}n_z$ and with
$H_{\mathrm{ph}}$, and the effective Hamiltonian is
\begin{equation}\label{eq:Hcold}
H_{\mathrm{cold}}=P_{\mathrm{cold}}\,e^{S_1}H_{\mathrm{LF}}e^{-S_1}\,P_{\mathrm{cold}},
\end{equation}
an operator on the finite-dimensional space $\mathcal{H}_{\mathrm{cold}}$. Its
Gibbs state is
\begin{equation}\label{eq:coldstate}
\langle\cdot\rangle_{\mathrm{cold}}
=\frac{\Tr_{\mathcal{H}_{\mathrm{cold}}}\bigl(\cdot\,e^{-\beta H_{\mathrm{cold}}}\bigr)}
{\Tr_{\mathcal{H}_{\mathrm{cold}}}\bigl(e^{-\beta H_{\mathrm{cold}}}\bigr)} .
\end{equation}
For a site-wise angle vector $\theta=(\theta_z)_{z\in\Lambda}\in\mathbb{R}^{\Lambda}$
let $R(\theta)=\exp\bigl(\sum_z\theta_zS_z\bigr)$ and
$\delta_e=\theta_u-\theta_v$ for $e=(u\leftarrow v)$. Since $R(\theta)$ commutes
with $P_{\mathrm{cold}}$, $H_{\mathrm{ph}}$ and $D_0$, and since
$[S_u,c^\dagger_{u\sigma}c_{v\sigma}]=+\tfrac{1}{2}c^\dagger_{u\sigma}c_{v\sigma}$,
the rotated effective Hamiltonian
$H_{\mathrm{cold}}(\theta)=R(\theta)H_{\mathrm{cold}}R(\theta)^{-1}$
equals $P_{\mathrm{cold}}e^{S_1(\theta)}H_{\mathrm{LF}}(\theta)e^{-S_1(\theta)}P_{\mathrm{cold}}$,
where every hop across $e$ from $v$ to $u$ carries the factor
$e^{+\delta_e/2}$ and every hop from $u$ to $v$ the factor $e^{-\delta_e/2}$.
The rotation multiplies $\Dop_{x_0}^\dagger$ by $e^{+\theta_{x_0}}$ and
$\Dop_{x_0}$ by $e^{-\theta_{x_0}}$, because $\Dop_{x_0}^\dagger$ creates and
$\Dop_{x_0}$ annihilates two electrons at $x_0$.

\subsection{Main statements}\label{sec:statements}

\begin{theorem}[Cold-projected pairing ceiling with an $e^{-7U/(8\Omega)}$
narrowing]\label{thm:main}
Let $U>0$ and $t\ge0$ satisfy $x=t/U\le2^{-14}$, let $X=U/\Omega\ge1$ and let
$M=\lfloor X/64\rfloor$. Write $K_4(x)$ for the series \eqref{eq:K4} below and
\begin{equation}\label{eq:A}
A(\beta)=8\beta U K_4(x)\,e^{-\frac{7}{8}X}>0 .
\end{equation}
Then the following hold on $\mathcal{H}_{\mathrm{cold}}$ for every side
$L\ge4$ and every real $\mu$.
\begin{enumerate}
\item\label{it:bch} The series
$H_{\mathrm{cold}}(\theta)=\sum_{m\ge0}\frac{1}{m!}P_{\mathrm{cold}}\ad_{S_1(\theta)}^{m}\bigl(H_{\mathrm{LF}}(\theta)\bigr)P_{\mathrm{cold}}$
converges in operator norm on $\mathcal{H}_{\mathrm{cold}}$ whenever
$|\theta_u-\theta_v|\le4\ln2$ on every edge.
\item\label{it:decay} For distinct sites $x_0\ne y_0$ with $d=d(x_0,y_0)\ge1$,
\begin{equation}\label{eq:decay}
\bigl|\langle\Dop_{x_0}^\dagger\Dop_{y_0}\rangle_{\mathrm{cold}}\bigr|
\le(d+1)^{-\xi_{\mathrm{cold}}(\beta)},\qquad
\xi_{\mathrm{cold}}(\beta)=
\begin{cases}
4-16A(\beta), & A(\beta)<\tfrac18,\\[2pt]
\dfrac{1}{4A(\beta)}=\dfrac{1}{32\beta UK_4(x)e^{-\frac{7}{8}X}}, & A(\beta)\ge\tfrac18 .
\end{cases}
\end{equation}
\item\label{it:ceiling} If
\begin{equation}\label{eq:ceiling}
\tau>2^{26}\,\frac{t^{2}}{U}\,e^{-\frac{7}{8}X},
\end{equation}
then $A(\beta)<1/8$, hence $\xi_{\mathrm{cold}}(\beta)>2$ and the equal-time
polaron-pair correlations are uniformly summable in the row,
\begin{equation}\label{eq:summ}
\sum_{y_0\in\Lambda}\bigl|\langle\Dop_{x_0}^\dagger\Dop_{y_0}\rangle_{\mathrm{cold}}\bigr|
\le1+\sum_{r\ge1}4r\,(r+1)^{-\xi_{\mathrm{cold}}(\beta)}<\infty ,
\end{equation}
the bound being independent of $L\ge4$ and of $\mu\in\mathbb{R}$. The exponent
$\xi_{\mathrm{cold}}(\beta)>2$ and the prefactor $A_0=1$ of \eqref{eq:decay} do
not depend on $L$ or $\mu$ at fixed $X\ge1$, which is the quantifier pattern of
Definition~\ref{def:tpair}. Consequently
$k_{\mathrm{B}}T^{\mathrm{cold}}_{\mathrm{pair}}\le2^{26}(t^{2}/U)e^{-7U/(8\Omega)}$
in the notation of Definition~\ref{def:tpair}.
\end{enumerate}
\end{theorem}

\begin{corollary}[The original pair field in the cold state]\label{cor:orig}
Under the hypotheses of Theorem~\ref{thm:main} let
\begin{equation}\label{eq:origpair}
\tilde\Dop^{\mathrm{orig}}_z
=P_{\mathrm{cold}}e^{S_1}U_{\mathrm{LF}}\Dop_zU_{\mathrm{LF}}^\dagger e^{-S_1}P_{\mathrm{cold}}
=P_{\mathrm{cold}}e^{S_1}\Dop_zD_z(+2\alpha)e^{-S_1}P_{\mathrm{cold}}
\end{equation}
be the cold-projected original electronic pair field. Then for all
$d=d(x_0,y_0)\ge1$,
\begin{equation}\label{eq:corbound}
\bigl|\bigl\langle(\tilde\Dop^{\mathrm{orig}}_{x_0})^\dagger\tilde\Dop^{\mathrm{orig}}_{y_0}\bigr\rangle_{\mathrm{cold}}\bigr|
\le4\,(d+1)^{-\xi_{\mathrm{cold}}(\beta)},
\end{equation}
so the ceiling \eqref{eq:ceiling} of Theorem~\ref{thm:main} also gives uniform
two-dimensional summability of the original electronic pair correlations inside
$\mathcal{H}_{\mathrm{cold}}$.
\end{corollary}

The exponent of \eqref{eq:decay} is the one produced by the optimisation of the
rotation strength over $\kappa\in(0,4]$ in Lemma~\ref{lem:holder} below, and the
constant $2^{26}$ is the round number that dominates the certified value of the
series. The decay is algebraic at every temperature, but the algebraic rate is
only larger than the summability threshold $2$ above the ceiling
\eqref{eq:ceiling}. The narrowing factor $e^{-7U/(8\Omega)}$ is what turns the
strong-coupling ceiling $2^{26}t^{2}/U$ into a bound below $2^{27}t^{2}/U$ for
every $X\ge1$, below $6000\,t^{2}/U$ once
$X>\frac{8}{7}\ln(2^{26}/6000)\approx10.65$, and, at fixed $t$ and $U$, into a
bound that decreases with $1/\Omega$ faster than any power. The reduction to the full Gibbs state is not made, and
Proposition~\ref{prop:thermal} explains why the present route cannot make it.

Table~\ref{tab:claims} collects the four statements of the note with their status
and with the mechanism that carries each one.

\begin{table}[t]
\centering
\small
\begin{tabular}{@{}L{0.30\textwidth}L{0.40\textwidth}L{0.22\textwidth}@{}}
\hline
Statement & Content & Status \\
\hline
Theorem~\ref{thm:main}, Corollary~\ref{cor:orig} &
Algebraic decay in the cold-projected effective model with prefactor
$2^{26}$ in the ceiling \eqref{eq:ceiling}, for the polaron pairs and for the
original pair field &
proved, with a certified replay of every constant \\
Proposition~\ref{prop:floor} &
Neutral excursions on two sites force a
$\tfrac{1}{400}(t^{2}/U)\sinh^{2}(\delta/2)$ free-energy gain, and the
Franck--Condon resolvent obeys a $\tfrac{1}{2U}$ floor, so a direct microscopic
rotation gains nothing below a constant multiple of $t^{2}/U$ &
proved, with a certified replay and a two-site numerical cross-check \\
Proposition~\ref{prop:thermal} &
Thermal pair hopping is bounded by $t^{2}/(eU^{2})$ in the crossover regime,
and the cold manifold has Gibbs weight at most
$\exp(-|\Lambda|e^{-\beta\Omega(M+1)})$ &
proved, with a certified peak identity and a numerical crossover check \\
Lifting Theorem~\ref{thm:main} to the full Gibbs state &
Admitting all phonon occupations $n^{\mathrm{ph}}_z>M$ without paying the
neutral-loop floor of Proposition~\ref{prop:floor} &
open \\
\hline
\end{tabular}
\caption{The statements of this note, their status, and their mechanism.}
\label{tab:claims}
\end{table}

\section{Proof of Theorem~\ref{thm:main} and Corollary~\ref{cor:orig}}\label{sec:proof}

The proof has the following shape. The rotation acts on the conjugated Hamiltonian
through the phases attached to hops, so $H_{\mathrm{cold}}(\theta)$ is expanded as
a series in $\theta$-dependent single-edge operators, and the twist part of that
series vanishes identically for every monomial that conserves the electronic
charge site by site. The second derivative of the remaining monomials with respect
to the rotation parameter carries the square of a net charge transfer, which for a
monomial that starts and ends with no singly occupied site is at least two on at
least two distinct sites. Each unit of net charge transfer at a site pairs with a
phonon displacement of magnitude $\alpha$ at that site, and the matrix element of
such a coherent displacement between phonon occupations at most $M+1$ and at most
$M$ is exponentially small in $X$. What remains is a counting problem for
connected chains of edges, which is Lemma~\ref{lem:clusters}, followed by a
one-dimensional summation along shells (Lemma~\ref{lem:shell}) and the
optimisation of the rotation strength (Lemma~\ref{lem:holder}).

\subsection{Norms of the edge generators and the cancellation identity}

\begin{lemma}\label{lem:norms}
For every edge $e$ and every real $\theta$,
\begin{equation}\label{eq:s1norm}
\|S_{1,e}(\theta)\|_{\mathrm{op}}\le2\sqrt{2}\cosh\!\bigl(2\ln2\bigr)x
=\frac{17\sqrt2}{4}\,x<7x \qquad\text{whenever } |\delta_{e}|=|\theta_u-\theta_v|\le4\ln2 ,
\end{equation}
the bound being uniform on all of $\mathcal{H}_\Lambda$. Moreover
$[S_1(\theta),D_0]=-V_{\mathrm{mix}}(\theta)$ identically, so $D_0+V_{\mathrm{mix}}(\theta)$
is the sum of an operator and its own Schrieffer--Wolff generator, and the series
\begin{equation}\label{eq:bch}
H_{\mathrm{cold}}(\theta)=\sum_{m\ge0}\frac{1}{m!}P_{\mathrm{cold}}\,
\ad_{S_1(\theta)}^{m}\bigl(H_{\mathrm{LF}}(\theta)\bigr)P_{\mathrm{cold}}
\end{equation}
converges in operator norm on $\mathcal{H}_{\mathrm{cold}}$ for every $\theta$
with $|\delta_e|\le4\ln2$ on every edge.
\end{lemma}

\begin{proof}
On an edge, $A_{e,\mathrm{mix}}$ is a matrix with entries of modulus $0$ or $1$
and $\|A_{e,\mathrm{mix}}\|_{\mathrm{op}}=\sqrt2$ by \eqref{eq:Anorms}; the
displacement operators are unitary; and the two signs of the twist contribute at
most $\cosh(2\ln2)=17/8$ in the worst case $|\delta_e|=4\ln2$, which gives the
first bound. The exact rational identities
$\|A_{e,\mathrm{mix}}\|_{\mathrm{op}}=\sqrt2$ and
$\cosh(2\ln2)=17/8$ are certified. For the cancellation, $S_1$ commutes with $N$
and hence with $\mu N$, and
$\bigl[\tfrac{U}{2}N_{\mathrm{odd}},P_{e,0}A_eP_{e,2}\bigr]
=-UP_{e,0}A_eP_{e,2}$ while
$\bigl[\tfrac{U}{2}N_{\mathrm{odd}},P_{e,2}A_eP_{e,0}\bigr]=+UP_{e,2}A_eP_{e,0}$,
so $[S_{1,e},\tfrac{U}{2}N_{\mathrm{odd}}]=-V_{e,\mathrm{mix}}$ for every $e$,
and summing over edges gives the identity. For convergence, $S_1(\theta)$ is
bounded on $\mathcal{H}_\Lambda$ and shifts phonon numbers by bounded
displacements $D_z(\pm\alpha)$, so for every $\psi\in\mathcal{H}_{\mathrm{cold}}$
and $\lambda\in\mathbb{C}$ the vector $e^{-\lambda S_1(\theta)}\psi$ lies in the
core $\mathcal{D}(H_{\mathrm{ph}})$. Hence
$\lambda\mapsto P_{\mathrm{cold}}e^{\lambda S_1(\theta)}H_{\mathrm{LF}}(\theta)
e^{-\lambda S_1(\theta)}P_{\mathrm{cold}}$ is an entire
$\mathcal{B}(\mathcal{H}_{\mathrm{cold}})$-valued function on finite $\Lambda$.
Its Taylor series at $\lambda=0$ is \eqref{eq:bch}, which converges at
$\lambda=1$ by the uniform tree/cluster bounds of Lemma~\ref{lem:clusters}: the
$m$-th term is bounded by a polynomial in $m$ times $r^{m}$ with
$r=56x\le7/2048$, uniformly in $L$ and $\mu$.
\end{proof}

Note that the first inequality of \eqref{eq:s1norm} is a bound on a single edge
generator, not on $S_1$; the norm of $S_1$ itself grows with the number of edges,
and it is the connectedness of the chains in Lemma~\ref{lem:clusters} that
prevents that growth from reaching the final bound.

\subsection{Charge-diagonal monomials carry no twist}

\begin{lemma}\label{lem:cancel}
On $\mathcal{H}_{\mathrm{cold}}$ the following hold for every $\theta$.
\begin{enumerate}
\item\label{it:ph} $P_{\mathrm{cold}}[S_1(\theta),H_{\mathrm{ph}}]P_{\mathrm{cold}}=0$
identically.
\item\label{it:diag} $P_{\mathrm{cold}}[S_1(\theta),V_{\mathrm{diag}}(\theta)]P_{\mathrm{cold}}=0$
identically.
\item\label{it:edges} For distinct edges $e_1\ne e_0$,
$P_{\mathrm{cold}}[S_{1,e_1}(\theta),V_{e_0,\mathrm{mix}}(\theta)]P_{\mathrm{cold}}=0$.
\item\label{it:monomial} Every monomial
$\mathcal{M}_{\gamma,\sigma}(s\theta)$ obtained by expanding
$H_{\mathrm{cold}}(s\theta)$ along a chain of edges and choosing a direction for
each hop has the exact form
\begin{equation}\label{eq:monomial}
\mathcal{M}_{\gamma,\sigma}(s\theta)=e^{s\Phi(\gamma,\sigma)}\mathcal{M}_{\gamma,\sigma}(0),
\qquad
\Phi(\gamma,\sigma)=\frac{1}{2}\sum_{j}\sigma_j\delta_{e_j}
=\frac{1}{2}\sum_{z\in\Lambda}\theta_z\,\Delta n_z(\gamma,\sigma),
\end{equation}
with $\Delta n_z(\gamma,\sigma)\in\mathbb{Z}$ the net electronic charge change at
site $z$. If $\Delta n_z(\gamma,\sigma)=0$ for every $z$, then
$\Phi(\gamma,\sigma)=0$ and $\frac{d^{2}}{ds^{2}}\mathcal{M}_{\gamma,\sigma}(s\theta)$
vanishes identically.
\end{enumerate}
\end{lemma}

\begin{proof}
Items 1 and 2 follow from the site-occupation structure. Every term of
$S_1(\theta)$ maps $\Nodd=0$ to $\Nodd=2$ and vice versa, while $H_{\mathrm{ph}}$
preserves $\Nodd$, so the two orders of the product map $\Nodd=0$ into
$\Nodd\ge2$, which the cold projection kills. The operator
$V_{\mathrm{diag}}(\theta)$ preserves $\Nodd$: it annihilates the states with both
sites even (there $A_{e,\mathrm{diag}}=0$) and preserves $\Nodd=2$, so both
$S_{1}V_{\mathrm{diag}}$ and $V_{\mathrm{diag}}S_1$ map $\Nodd=0$ to
$\Nodd\ge2$. For item 3, the edge $e_0$ creates two odd sites and the edge
$e_1\ne e_0$ can flip at most one of them back, so the intermediate state still
has at least two odd sites and is annihilated by $P_{\mathrm{cold}}$. Item 4 is a
computation: each hop across $e=(u\leftarrow v)$ contributes $e^{+s\delta_e/2}$ in
the direction that increases the charge at $u$, and the product over the chain
collects the sum in the first identity; the second identity is the same statement
written in terms of the charge differences at the sites. A monomial that conserves
the charge at every site therefore has no $s$ dependence, and its second
derivative vanishes.
\end{proof}

Two consequences are used below. First, the terms with $m=0$ and $m=1$ of
\eqref{eq:bch} collapse: $D_0+V_{\mathrm{mix}}(\theta)+[S_1(\theta),D_0]=D_0$, so
$m=0$ contributes $P_{\mathrm{cold}}(D_0+H_{\mathrm{ph}}+V_{\mathrm{diag}}(\theta))P_{\mathrm{cold}}
=P_{\mathrm{cold}}(D_0+H_{\mathrm{ph}})P_{\mathrm{cold}}$ by item 2, and $m=1$
reduces to the three commutators of item 1 to 3. Second, for $m\ge2$ the
$(m+1)$-st commutator of $D_0$ combines with the $m$-th commutator of
$V_{\mathrm{mix}}(\theta)$, because
$\frac{1}{(m+1)!}\ad^{m+1}_{S_1}(D_0)+\frac{1}{m!}\ad^{m}_{S_1}(V_{\mathrm{mix}})
=\frac{m}{(m+1)!}\ad_{S_1}^{m}(V_{\mathrm{mix}})$ by
Lemma~\ref{lem:norms}, and $0<\frac{m}{(m+1)!}<\frac{1}{m!}$. The terms of the
expansion therefore organise into two families of connected chains,
\begin{align}
\text{Family I:}\quad &\gamma=(e_1,\dots,e_m;e_0),\qquad
\frac{1}{m!}P_{\mathrm{cold}}\ad_{S_{1,e_1}}\cdots\ad_{S_{1,e_m}}
\bigl(V_{e_0,\mathrm{mix}}+V_{e_0,\mathrm{diag}}\bigr)P_{\mathrm{cold}},
\qquad m\ge1,
\label{eq:famI}\\
\text{Family II:}\quad &\gamma=(e_1,\dots,e_m),\qquad
\frac{1}{m!}P_{\mathrm{cold}}\ad_{S_{1,e_1}}\cdots\ad_{S_{1,e_{m-1}}}
\bigl[S_{1,e_m},H_{\mathrm{ph}}\bigr]P_{\mathrm{cold}},\qquad m\ge2,
\label{eq:famII}
\end{align}
where in Family I the case $m=1$ absorbs the factor $\tfrac{1}{2}$ of the
commutator with $V_{e_0,\mathrm{mix}}$ coming from the collapse above.

\subsection{The Hermitian twist remainder}

Following the complex-rotation method, define
\begin{equation}\label{eq:C}
C(\theta)=\frac12\bigl(H_{\mathrm{cold}}(\theta)+H_{\mathrm{cold}}(-\theta)\bigr)-H_{\mathrm{cold}}(0)
=\frac12\int_0^1(1-s)\Bigl(\frac{d^2}{ds^2}H_{\mathrm{cold}}(s\theta)+\frac{d^2}{ds^2}H_{\mathrm{cold}}(-s\theta)\Bigr)ds ,
\end{equation}
which is Hermitian and satisfies $C(0)=0$. The first variation of
$H_{\mathrm{cold}}$ is a sum of monomials linear in $s$, and by
Lemma~\ref{lem:cancel}(4) the only monomials that survive the second derivative
are those with a net charge transfer of modulus at least two at two distinct
sites. The following lemma is the analytical heart of the estimate.

\begin{lemma}[Coherent-state bound]\label{lem:bf}
For every $X=U/\Omega>0$, $M=\lfloor X/64\rfloor$, $\alpha=\sqrt{X/2}$ and every
integer $|s|\ge2$,
\begin{equation}\label{eq:bf}
\bigl\|P_{\le M+1}D_z(s\alpha)P_{\le M}\bigr\|_{\mathrm{op}}\le\sqrt2\,e^{-\frac{7}{16}X},
\end{equation}
where $P_{\le n}$ is the projector onto phonon occupations $0,\dots,n$ at the
site $z$.
\end{lemma}

\begin{proof}
Let $\beta=s\alpha$, so $|\beta|\ge\sqrt{2X}$. If $X<64$ then $M=0$ and the only
matrix elements are $\langle0|D(\beta)|0\rangle=e^{-\beta^{2}/2}$ and
$\langle1|D(\beta)|0\rangle=\beta e^{-\beta^{2}/2}$, so
$\|P_{\le1}D(\beta)P_{\le0}\|_{\mathrm{op}}=\sqrt{1+\beta^{2}}\,e^{-\beta^{2}/2}
\le\sqrt2\,e^{-\beta^{2}/4}\le\sqrt2\,e^{-X/2}$, using $1+u\le2e^{u/2}$ for $u\ge0$.
If $X\ge64$ then $M\ge1$ and the Bargmann--Fock generating identity
\begin{equation}\label{eq:bfgen}
\sum_{n,m\ge0}\frac{w^{n}z^{m}}{\sqrt{n!m!}}\langle n|D(\beta)|m\rangle
=\exp\Bigl(-\frac{\beta^{2}}{2}+\beta w-\beta z+wz\Bigr)
\end{equation}
together with Cauchy's formula on $|w|=|z|=r:=\sqrt{M+1}$ gives, for
$0\le n\le M+1$ and $0\le m\le M$,
\begin{equation}\label{eq:bfentry}
\bigl|\langle n|D(\beta)|m\rangle\bigr|
\le\frac{\sqrt{n!m!}}{r^{n+m}}\exp\Bigl(-\frac{\beta^{2}}{2}+2|\beta|r+r^{2}\Bigr)
\le\exp\Bigl(-\frac{\beta^{2}}{2}+2|\beta|r+r^{2}\Bigr),
\end{equation}
because $n!\le(M+1)^{n}=r^{2n}$ and $m!\le(M+1)^{m}=r^{2m}$. Since $M\le X/64$ and
$1\le X/64$ we have $r^{2}\le X/32$, and hence
\begin{equation}\label{eq:bfexp}
-\frac{\beta^{2}}{2}+2|\beta|r+r^{2}
\le-\frac{s^{2}X}{4}+2|s|\sqrt{\frac X2}\sqrt{\frac X{32}}+\frac X{64}+1
=1-\Bigl(s^{2}-|s|-\frac1{16}\Bigr)\frac X4\le1-\frac{31}{64}X
\end{equation}
for $|s|\ge2$. Frobenius domination on the $(M+2)\times(M+1)$ matrix, together
with $r^{2}\le X/32$, gives
\begin{equation}\label{eq:bfdir}
\|P_{\le M+1}D(\beta)P_{\le M}\|_{\mathrm{op}}
\le\sqrt{(M+2)(M+1)}\,e^{1-\frac{31}{64}X}
\le\Bigl(\frac X{64}+2\Bigr)e^{1-\frac{3}{64}X}e^{-\frac{7}{16}X}\le\sqrt2\,e^{-\frac{7}{16}X},
\end{equation}
because $X\mapsto(\frac X{64}+2)e^{1-\frac{3}{64}X}$ is strictly decreasing on
$X\ge0$ and equals $3e^{-2}<\sqrt2$ at $X=64$. The final comparison at
$X=64$ and the sign of the derivative $\frac1{128}-\frac3{64}$ are certified.
\end{proof}

Applied at two distinct sites, \eqref{eq:bf} gives the factor
$(\sqrt2e^{-7X/16})^{2}=2e^{-7X/8}$, and at all other sites the compression
$P_{\le M}D_z(\cdot)P_{\le M}$ has norm at most one. This factor of
$2e^{-7X/8}$ pays for the pair of displacement operators created by a net
transferred charge of modulus two at two distinct sites in
Lemma~\ref{lem:cancel}(4).

Rescaling by $U$: for an integer $m\ge1$ and $X\ge1$ the norm on
$\mathcal{H}_{\le M}$ of the Weyl pull of the oscillator energy,
\begin{equation}\label{eq:phpull}
P_{\le M}D_z(s_z\alpha)^\dagger\Omega a_z^\dagger a_zD_z(s_z\alpha)P_{\le M}
=P_{\le M}\Omega\bigl(a_z^\dagger+s_z\alpha\bigr)\bigl(a_z+s_z\alpha\bigr)P_{\le M},
\end{equation}
with $|s_z|\le m$ is at most $\Omega(\sqrt{M+1}+m\sqrt{X/2})^{2}\le U(m+2)^{2}$,
because $M+1\le X/64+1\le\frac{65}{64}X<4X$ for $X\ge1$ and $\Omega=X^{-1}U$.
This is the factor that appears in Family II of \eqref{eq:famII}.

\subsection{Connected chains and the counting lemma}

For a site $x_0$ and a second site $y_0$ at distance $d=d(x_0,y_0)\ge1$ set
\begin{equation}\label{eq:theta}
\theta(z)=-\kappa\max\Bigl(0,\ln\frac{d+1}{d(x_0,z)+1}\Bigr),\qquad 0<\kappa\le4 .
\end{equation}
On every edge $e=\{u,v\}$ one has
$|\delta_e|=|\theta(u)-\theta(v)|\le\kappa\ln2\le4\ln2$, since the ratio of the two
arguments of the logarithm is at most $2$ on adjacent sites and $\kappa\le4$, the
bound $4\ln2=2.7726$ being attained only at $\kappa=4$ and then only on an edge
across which the profile decays. Then $|\theta(x_0)-\theta(y_0)|=\kappa\ln(d+1)$, which is the gain used
in Lemma~\ref{lem:holder}, and $\theta(y_0)=0$.

\begin{lemma}[Chains and the operator norm of $C(\theta)$]\label{lem:clusters}
Under \eqref{eq:theta}, for every $x=t/U\le2^{-14}$ and every $X\ge1$,
\begin{equation}\label{eq:K4}
\bigl\|C(\theta)\bigr\|_{\mathrm{op}}
\le U\,K_4(x)\,e^{-\frac78X}\sum_{e\in E(\Lambda)}\delta_e^{2}
\end{equation}
with
\begin{equation}\label{eq:K4def}
K_4(x)=\sum_{m\ge1}2(m+1)^{5}(12x)\,r^{m}+\sum_{m\ge2}m^{4}(m+2)^{2}r^{m}
\le\frac{2^{26}}{64}x^{2}=2^{20}x^{2},\qquad r=56x\le\frac7{2048}<\frac1{256} .
\end{equation}
\end{lemma}

\begin{proof}
Fix one monomial. By Lemma~\ref{lem:cancel}(4), applied to the pair of monomials
$\mathcal{M}_{\gamma,\sigma}(\pm s\theta)$ which occur in \eqref{eq:C}, the
second derivative is $\Phi(\gamma,\sigma)^{2}$ times the monomial, and by
Cauchy--Schwarz
\begin{equation}\label{eq:phi2}
\Phi(\gamma,\sigma)^{2}\le\frac{m+1}{4}\sum_{a=0}^{m}\delta_{e_a}^{2}
\ \ \text{(Family I)},\qquad
\Phi(\gamma,\sigma)^{2}\le\frac{m}{4}\sum_{a=1}^{m}\delta_{e_a}^{2}
\ \ \text{(Family II)} .
\end{equation}
For the remaining factors, summing over the two direction choices of bond $e_j$
at fixed $s\in[0,1]$ costs at most
$\|S_{1,e_j}(s\theta)\|_{\mathrm{op}}\le2\sqrt2\cosh(2\ln2)x=\frac{17\sqrt2}{4}x<7x$
by Lemma~\ref{lem:norms}, and each commutator $\ad_{S_{1,e_j}(s\theta)}$
multiplies on the left and on the right, which costs $2\cdot7x=14x$. The pair
$V_{e_0,\mathrm{mix}}+V_{e_0,\mathrm{diag}}$ costs at most
$\|V_{e_0,\mathrm{mix}}\|_{\mathrm{op}}+\|V_{e_0,\mathrm{diag}}\|_{\mathrm{op}}
<(7+5)t=12Ux$, where $\|V_{e_0,\mathrm{mix}}\|_{\mathrm{op}}\le2\sqrt2\cosh(2\ln2)t
=\frac{17\sqrt2}{4}t<7t$ by the same computation as in Lemma~\ref{lem:norms}, the
two displacement operators having norm one, and
$\|V_{e_0,\mathrm{diag}}\|_{\mathrm{op}}\le2\cosh(2\ln2)t=\frac{17}{4}t<5t$
because $\|A_{e,\mathrm{diag}}\|_{\mathrm{op}}=1$.
Lemma~\ref{lem:bf} contributes $2e^{-7X/8}$ at the two sites carrying the net
charge, and the weight $\int_0^{1}(1-s)ds=\frac12$ of \eqref{eq:C} combines with
it into the single factor $e^{-7X/8}$.

We now bound the number of non-vanishing nested-commutator edge sequences
containing a fixed edge $e$. A nested commutator vanishes unless each outer
operator shares at least one vertex with the support of the inner commutator on
which it acts, so the edges of a non-vanishing sequence form a connected set.
Fix the edge $e$ and suppose that $e=e_a$ for some slot $a$.

\emph{Where the innermost edge can be.} If the sequence contains $e$ and is
connected, then in Family I its innermost edge $e_0$ lies within graph distance
at most $m$ of $e$, and in Family II its innermost edge $e_m$ lies within graph
distance at most $m-1$ of $e$. Since the union of the two balls of radius $R$
around the endpoints of $e$ has at most $2(R+1)^{2}$ vertices and the torus has
degree four, there are at most $8(m+1)^{2}$ admissible choices for $e_0$ in
Family I and at most $8m^{2}$ for $e_m$ in Family II.

\emph{Growing the cluster inside-out.} In Family II we enumerate the sequence
from the innermost edge $e_m$ backwards: $(e_m, e_{m-1}, \dots, e_1)$. Suppose
$e_{k+1},\dots,e_m$ have been placed, so that $r=m-k$ edges are already fixed.
For $\ad_{S_{1,e_k}}$ not to vanish on the inner commutator of
$\{e_{k+1},\dots,e_m\}$, the edge $e_k$ must share a vertex with the active
vertex set of $\{e_{k+1},\dots,e_m\}$, which spans at most $r+1$ vertices and
hence carries at most $4(r+1)$ incident edges. Running over
$k=m-1,m-2,\dots,1$, that is over $r=1,2,\dots,m-1$, gives at most
\begin{equation}\label{eq:countII}
\prod_{r=1}^{m-1}4(r+1)=4^{m-1}\,m!
\end{equation}
choices for $(e_{m-1},\dots,e_1)$ once $e_m$ is fixed. In Family I the same
inside-out enumeration starts from the innermost edge $e_0$ and places
$(e_m,e_{m-1},\dots,e_1)$ in reverse order: after $e_0$ and $r-1$ edges
$e_m,\dots,e_{m-r+2}$ have been placed ($r$ edges in total, $1\le r\le m$), the
active vertex set spans at most $r+1$ vertices, so the next edge $e_{m-r+1}$ has
at most $4(r+1)$ choices. Multiplying over $r=1,2,\dots,m$ gives at most
\begin{equation}\label{eq:countI}
\prod_{r=1}^{m}4(r+1)=4^{m}\,(m+1)!
\end{equation}
choices for $(e_m,\dots,e_1)$ once $e_0$ is fixed. Because the commutator nests
from the inside out, it is the suffix of already-placed inner edges that is
guaranteed to stay connected at every step.

\emph{The two series.} For Family I with $m\ge1$, summing over the $m+1$ possible
slots $a$ of $e$, using \eqref{eq:phi2}, the count $8(m+1)^{2}$ for $e_0$, the
count \eqref{eq:countI}, the factor $e^{-7X/8}$, the factor $12Ux$ and the factor
$(14x)^{m}/m!$ from the $m$ commutators, each $m$ contributes at most
\begin{align}
&e^{-\frac78X}\cdot(m+1)\cdot\frac{m+1}{4}\cdot8(m+1)^{2}\cdot\frac{4^{m}(m+1)!}{m!}\cdot(12Ux)(14x)^{m}\delta_e^{2}
\nonumber\\
&\qquad=U\,e^{-\frac78X}\,2(m+1)^{5}(12x)\,r^{m}\,\delta_e^{2},
\label{eq:famIterm}
\end{align}
where the powers of $m+1$ and of $4$ and $14$ were collected as
$(m+1)^{5}/4\cdot8=2(m+1)^{5}$ and $4^{m}14^{m}x^{m}=r^{m}$ with $r=56x$. For
Family II with $m\ge2$, the same computation with the slot sum $\sum_{a=1}^{m}$
giving $m\cdot\frac{m}{4}$, the count $8m^{2}$ for the innermost edge $e_m$, the count
\eqref{eq:countII}, the factor $2U(m+2)^{2}$ from \eqref{eq:phpull} and the factor
$(14x)^{m}/m!$ gives at most
\begin{equation}\label{eq:famIIterm}
e^{-\frac78X}\cdot m\cdot\frac m4\cdot8m^{2}\cdot\frac{4^{m-1}m!}{m!}\cdot2U(m+2)^{2}(14x)^{m}\delta_e^{2}
=U\,e^{-\frac78X}\,m^{4}(m+2)^{2}r^{m}\,\delta_e^{2} .
\end{equation}
Summing over $m$ and over the edges gives \eqref{eq:K4} with the series
\eqref{eq:K4def}. The numerical bound on the series is exact rational arithmetic:
summing the first twenty-nine terms of both series in exact fractions and bounding
the remaining tails by the exact geometric majorant gives
$64K_4(x)/x^{2}<5.562\times10^{7}<2^{26}\approx6.711\times10^{7}$ at the endpoint
$x=2^{-14}$, which is the only case that needs to be checked because $K_4(x)/x^{2}$
is increasing in $x$ and the series is increasing in $x$ term by term. Hence
$K_4(x)\le2^{20}x^{2}$ for $x\le2^{-14}$, and $r=56x\le56/16384=7/2048$.
\end{proof}

\begin{remark}\label{rem:noabs}
The bound \eqref{eq:K4} already contains the convergence assertion of
Lemma~\ref{lem:norms}. Each term of \eqref{eq:bch} is a sum over chains, each
chain contributes at most the corresponding term of \eqref{eq:famIterm} or
\eqref{eq:famIIterm} with $\delta_e^{2}$ replaced by $\max_e\delta_e^{2}\le(4\ln2)^{2}$,
and the resulting series in $m$ is dominated by a fixed polynomial in $m$ times
$r^{m}$ with $r=56x\le7/2048$, so it converges absolutely in operator norm on
$\mathcal{H}_{\mathrm{cold}}$. There is no need for a bound on $\|S_1\|$ uniform in
$L$; the polynomial growth of the number of chains is compensated by the factorial
weight $1/m!$ and by $r<1/256$.
\end{remark}

\subsection{The shell sum}

\begin{lemma}\label{lem:shell}
Under \eqref{eq:theta} and for every $L\ge4$,
\begin{equation}\label{eq:shell}
\sum_{e\in E(\Lambda)}\delta_e^{2}\le8\kappa^{2}\ln(d+1).
\end{equation}
\end{lemma}

\begin{proof}
Since $\theta$ vanishes at distance at least $d$ from $x_0$, only the edges
between shells $r-1$ and $r$ for $r=1,\dots,d$ contribute. On the torus the
covering map $\mathbb{Z}^{2}\to\Lambda$ is non-expansive, so the number of such
edges is at most its $\mathbb{Z}^{2}$ count $8r-4$. Since $0<\ln\frac{r+1}{r}\le
\frac1r$ for $r\ge1$, and since $\kappa\ln\frac{r+1}{r}$ is the decrement of
$\theta$ across an edge from shell $r-1$ to shell $r$,
\begin{equation}\label{eq:shell2}
\sum_{e}\delta_e^{2}\le\kappa^{2}\sum_{r=1}^{d}(8r-4)\Bigl(\ln\frac{r+1}{r}\Bigr)^{2}
\le8\kappa^{2}\sum_{r=1}^{d}\ln\frac{r+1}{r}=8\kappa^{2}\ln(d+1),
\end{equation}
the middle inequality using
$(8r-4)\ln^{2}\frac{r+1}{r}\le8r\ln^{2}\frac{r+1}{r}\le8\ln\frac{r+1}{r}$ for
$r\ge1$, which holds because $\ln\frac{r+1}{r}\le\frac1r\le1$. The last step telescopes.
\end{proof}

\subsection{Trotter, H\"older and the optimisation of \texorpdfstring{$\kappa$}{kappa}}

\begin{lemma}\label{lem:holder}
Let $A(\beta)=8\beta UK_4(x)e^{-7X/8}$ as in \eqref{eq:A}. Then for all distinct
$x_0,y_0$,
\begin{equation}\label{eq:holderchain}
\bigl|\langle\Dop_{x_0}^\dagger\Dop_{y_0}\rangle_{\mathrm{cold}}\bigr|
\le(d+1)^{-\kappa}e^{\beta\|C(\theta)\|_{\mathrm{op}}}
\le(d+1)^{-\kappa+A\kappa^{2}} .
\end{equation}
Optimising over $\kappa\in(0,4]$ gives \eqref{eq:decay}, and
$A(\beta)<1/8$ holds whenever $\tau>2^{26}(t^{2}/U)e^{-7X/8}$.
\end{lemma}

\begin{proof}
Write $H_{\mathrm{cold}}(\theta)=H_{\mathrm{cold}}+B(\theta)+C(\theta)$ with
$B(\theta)=\tfrac12(H_{\mathrm{cold}}(\theta)-H_{\mathrm{cold}}(-\theta))$
anti-Hermitian, so that $B(\theta)^\dagger=-B(\theta)$ and $C(\theta)$ Hermitian.
Since $R(\theta)$ commutes with $P_{\mathrm{cold}}$ and with the trace on
$\mathcal{H}_{\mathrm{cold}}$, cyclicity gives
$\Tr\bigl(\Dop_{x_0}^\dagger\Dop_{y_0}e^{-\beta H_{\mathrm{cold}}}\bigr)
=\Tr\bigl(R(\theta)\Dop_{x_0}^\dagger\Dop_{y_0}R(\theta)^{-1}e^{-\beta H_{\mathrm{cold}}(\theta)}\bigr)$,
and $R(\theta)\Dop_{x_0}^\dagger\Dop_{y_0}R(\theta)^{-1}=e^{\theta(x_0)-\theta(y_0)}\Dop^\dagger_{x_0}\Dop_{y_0}
=(d+1)^{-\kappa}\Dop_{x_0}^\dagger\Dop_{y_0}$ because the pair field carries charge
two. Trotter's product formula \cite{trotter1959} for
$e^{-\beta(H_{\mathrm{cold}}+B+C)}$,
$\|e^{-(\beta/n)B(\theta)}\|_{\mathrm{op}}=1$ for the anti-Hermitian factor, and
H\"older's inequality for Schatten norms with the Peierls--Bogoliubov inequality
for the Hermitian factor $H_{\mathrm{cold}}+C(\theta)$ bound the ratio of traces
by $e^{\beta\|C(\theta)\|_{\mathrm{op}}}$, so
$\bigl|\langle\Dop^\dagger_{x_0}\Dop_{y_0}\rangle_{\mathrm{cold}}\bigr|
\le\|\Dop^\dagger_{x_0}\Dop_{y_0}\|_{\mathrm{op}}(d+1)^{-\kappa}
e^{\beta\|C(\theta)\|_{\mathrm{op}}}\le(d+1)^{-\kappa}e^{\beta\|C(\theta)\|_{\mathrm{op}}}$,
since the pair field has norm one. Inserting \eqref{eq:K4} and \eqref{eq:shell}
gives $\beta\|C(\theta)\|_{\mathrm{op}}\le A\kappa^{2}\ln(d+1)$ with $A=8\beta UK_4(x)e^{-7X/8}$,
so $e^{\beta\|C(\theta)\|_{\mathrm{op}}}\le(d+1)^{A\kappa^{2}}$.
Maximising $\kappa-A\kappa^{2}$ over $\kappa\in(0,4]$: if $A<1/8$ the maximum over
$\kappa\in(0,4]$ is at $\kappa=4$ and equals $4-16A$, so the left side of
\eqref{eq:holderchain} is at most $(d+1)^{-(4-16A)}$ with $4-16A>2$; if $A\ge1/8$
the unconstrained maximum is at $\kappa=1/(2A)\le4$ and equals $1/(4A)$. Finally,
$A(\beta)<1/8$ is equivalent to $\tau>64UK_4(x)e^{-7X/8}$, which follows from
$\tau>2^{26}(t^{2}/U)e^{-7X/8}$ because $64K_4(x)<2^{26}x^{2}$ by \eqref{eq:K4def}.
\end{proof}

\subsection{End of the proofs}

\begin{proof}[Proof of Theorem~\ref{thm:main}]
Item~\ref{it:bch} is Lemma~\ref{lem:norms}. Item~\ref{it:decay} is
Lemma~\ref{lem:holder} with \eqref{eq:decay} as stated. For
item~\ref{it:ceiling}, \eqref{eq:ceiling} gives $A(\beta)<1/8$, hence
$\xi_{\mathrm{cold}}(\beta)=4-16A(\beta)>2$ and the decay
\eqref{eq:decay} is summable. The number of sites at distance exactly $r$ from
$x_0$ is at most $4r$ on $\Lambda$ for $L\ge4$, so
$\sum_{y_0}\bigl|\langle\Dop^\dagger_{x_0}\Dop_{y_0}\rangle_{\mathrm{cold}}\bigr|
\le1+\sum_{r\ge1}4r(r+1)^{-\xi_{\mathrm{cold}}(\beta)}$, which is finite and
independent of $L$ and $\mu$ because the exponent bound does not involve $\mu$ and
the shell counts are torus-independent. The prefactor is $A_0=1$ and the exponent
$\xi_{\mathrm{cold}}(\beta)>2$, neither depending on $L$ or $\mu$ at fixed
$X\ge1$, so the property required by Definition~\ref{def:tpair} holds at every
$\tau$ above the right side of \eqref{eq:ceiling}, and the infimum in
Definition~\ref{def:tpair} is therefore at most that right side.
\end{proof}

\begin{proof}[Proof of Corollary~\ref{cor:orig}]
Since $U_{\mathrm{LF}}\Dop_zU_{\mathrm{LF}}^\dagger=\Dop_zD_z(+2\alpha)$ and
$U_{\mathrm{LF}}\Dop_z^\dagger U_{\mathrm{LF}}^\dagger=\Dop_z^\dagger D_z(-2\alpha)$,
Lemma~\ref{lem:norms} gives
$\tilde\Dop^{\mathrm{orig}}_z=P_{\mathrm{cold}}\bigl(\sum_{m\ge0}\frac1{m!}\ad_{S_1}^{m}
(\Dop_zD_z(+2\alpha))\bigr)P_{\mathrm{cold}}$. The rotation is an algebra
automorphism commuting with $P_{\mathrm{cold}}$ and acts on the pair field with
the displacement by $e^{\pm\theta(z)}$, because
$R(\theta)D_z(\pm2\alpha)R(\theta)^{-1}=D_z(\pm2\alpha)$, so
\begin{equation}\label{eq:corrot}
R(\theta)(\tilde\Dop^{\mathrm{orig}}_{x_0})^\dagger R(\theta)^{-1}
=(d+1)^{-\kappa}P_{\mathrm{cold}}\Bigl(\sum_{m\ge0}\frac1{m!}\ad_{S_1(\theta)}^{m}
\bigl(\Dop_{x_0}^\dagger D_{x_0}(-2\alpha)\bigr)\Bigr)P_{\mathrm{cold}} ,
\end{equation}
and the same at $y_0$ with $e^{-\theta(y_0)}=1$. In the $m$-th commutator the
$j$-th step chooses an edge touching an active set of at most $j+1$ vertices, at
most $4(j+1)$ choices, and contributes at most $2\|S_{1,e_j}(\theta)\|_{\mathrm{op}}<14x$;
the pair field and the displacement are of norm one. Dividing the product
$\prod_{j=1}^{m}4(j+1)=4^{m}(m+1)!$ by $m!$ gives $(m+1)(56x)^{m}=(m+1)r^{m}$, and
$\sum_{m\ge0}(m+1)r^{m}=(1-r)^{-2}<2$ for $r\le7/2048$, which is certified. The
same bound at $y_0$ gives
$\|R(\theta)(\tilde\Dop^{\mathrm{orig}}_{x_0})^\dagger\tilde\Dop^{\mathrm{orig}}_{y_0}R(\theta)^{-1}\|_{\mathrm{op}}
\le4(d+1)^{-\kappa}$, after which Lemma~\ref{lem:holder} applies verbatim.
\end{proof}

\section{Obstruction I: the neutral-loop floor for a direct microscopic rotation}\label{sec:floor}

The complex-rotation method can be applied to the Holstein Hamiltonian itself,
without a Lang--Firsov transformation, without a Schrieffer--Wolff truncation and
without the cold projection. One rotates the electronic charge by $\theta$, writes
the trace in the twisted Hamiltonian, separates the rotation into the Hermitian
twist, which multiplies each bond hopping by $\cosh(\delta_e/2)$, and the
anti-Hermitian part, whose norm is one, and then bounds the gain of the twist by
the free energy of the symmetrised Hamiltonian. On a bond, the twist multiplies the
hopping by $\cosh(\delta_e/2)\ge1$, and at second order in the hopping the
symmetrised Hamiltonian competes the pair-transfer channel $|2,0\rangle,
|0,2\rangle\leftrightarrow|\mathrm{sing}\rangle$ against the polaronic energy. The
competition is what Proposition~\ref{prop:floor} quantifies on two sites. It is
exact for every phonon frequency, and it is the reason the exponential narrowing
of Theorem~\ref{thm:main} cannot be seen without the cold projection.

\begin{proposition}[Neutral-loop floor on two sites]\label{prop:floor}
Let $t\ge0$, $U>0$, $\Omega>0$, and let $H$ be the two-site Hamiltonian
\eqref{eq:H} restricted to $N=2$, $S_z=0$ at $\mu=0$. For $\delta\in\mathbb{R}$
let $R(\delta)=e^{(\delta/2)n_0}$ and define the Hermitian twist part and its free
energy by
\begin{equation}\label{eq:hsym}
\begin{aligned}
H_{\mathrm{sym}}(\delta)&=\frac12\bigl(R(\delta)HR(\delta)^{-1}+R(-\delta)HR(-\delta)^{-1}\bigr)
=H_{\mathrm{at}}+\cosh(\delta/2)H_{\mathrm{hop}},\\
F_{\mathrm{sym}}(\delta)&=-\frac1\beta\ln\Tr_{N=2,S_z=0}e^{-\beta H_{\mathrm{sym}}(\delta)} .
\end{aligned}
\end{equation}
\begin{enumerate}
\item\label{it:fc} For all $\Omega>0$ and $U>0$, with $X=U/\Omega$ and
$D_{\mathrm{rel}}=D_0(-\alpha)D_1(+\alpha)$ the relative displacement of
magnitude $\sqrt2\alpha=\sqrt X$,
\begin{equation}\label{eq:fc}
\sum_{n\ge0}\frac{\bigl|\langle n|D_{\mathrm{rel}}|0\rangle\bigr|^{2}}{U+n\Omega}
=\sum_{n\ge0}e^{-X}\frac{X^{n}}{n!}\frac{1}{U+n\Omega}\ \ge\ \frac{1}{2U} .
\end{equation}
\item\label{it:gain} For all $\Omega>0$, all $\beta U\ge4$, all $0<t/U\le1/16$
and all $\delta$ with $\eta=\cosh(\delta/2)\in[1,2]$,
\begin{equation}\label{eq:gain}
-\bigl(F_{\mathrm{sym}}(\delta)-F_{\mathrm{sym}}(0)\bigr)
=\frac1\beta\ln\frac{\Tr\bigl(e^{-\beta H_{\mathrm{sym}}(\delta)}\bigr)}
{\Tr\bigl(e^{-\beta H_{\mathrm{sym}}(0)}\bigr)}
\ \ge\ \frac{1}{400}\,\frac{t^{2}}{U}\,\sinh^{2}(\delta/2) .
\end{equation}
\end{enumerate}
\end{proposition}

Part~\ref{it:fc} is the zero-temperature statement in the adiabatic
representation. Part~\ref{it:gain} is the finite-temperature statement, whose
proof sacrifices numerical sharpness to obtain a uniform bound for all
$\Omega>0$ and $\beta U\ge4$. At $T\to0$ and $t/U\to0$, part
\ref{it:fc} gives the asymptotic coefficient
$2\sum_{n}e^{-X}X^{n}/(n!\,(U+n\Omega))\ge1/U$, that is $c_0\ge1$ in
$F_{\mathrm{sym}}(\delta)-F_{\mathrm{sym}}(0)\simeq-c_0(t^{2}/U)\sinh^{2}(\delta/2)$,
which is the coefficient reproduced by the numerical check described below. In the
zero-displacement reference model ($g=0$) on the four-dimensional pair-singlet
subspace of an attractive Hubbard bond, the exact level splitting obeys
\begin{equation}\label{eq:hubbard}
E_0(t)-E_0\bigl(t\cosh(\delta/2)\bigr)\ge\frac{4t^{2}\sinh^{2}(\delta/2)}
{\sqrt{U^{2}+16t^{2}\cosh^{2}(\delta/2)}} ,
\end{equation}
which is certified as an algebraic identity in exact rational arithmetic.

\begin{proof}[Proof of Proposition~\ref{prop:floor}, part \ref{it:fc}]
Let $N$ be Poisson with parameter $X$, so that $\mathbb{P}(N=n)=e^{-X}X^{n}/n!$ and
$\mathbb{E}[N]=X=U/\Omega$. The function $g(y)=(U+y\Omega)^{-1}$ is strictly
convex on $[0,\infty)$, and Jensen's inequality gives
$\mathbb{E}[g(N)]\ge g(\mathbb{E}N)=(U+U)^{-1}=\frac{1}{2U}$, at every
$\Omega>0$. The Poisson weights are those of the relative coherent state,
$\bigl|\langle n|D_{\mathrm{rel}}|0\rangle\bigr|^{2}=e^{-X}X^{n}/n!$, which is the
identity listed in \eqref{eq:fc}.
\end{proof}

\begin{proof}[Proof of Proposition~\ref{prop:floor}, part \ref{it:gain}]
On two sites with $N=2$ and $S_z=0$ the electronic states are
$|2,0\rangle$, $|0,2\rangle$, the singlet $|\mathrm{sing}\rangle$ and the triplet
$|\mathrm{trip}\rangle$. The centre-of-mass mode
$a_{\mathrm{cm}}=(a_0+a_1)/\sqrt2$ decouples, because it commutes with
$(n_0-1,n_1-1)$ and with $H_{\mathrm{hop}}$, so the ratio of traces in
\eqref{eq:gain} equals the same ratio for the relative mode
$a_{\mathrm{rel}}=(a_0-a_1)/\sqrt2$. The triplet is annihilated by
$H_{\mathrm{hop}}$ and has Hamiltonian $U\mathbf 1+\Omega n_{\mathrm{rel}}$, so
its Gibbs trace is $Z_{\mathrm{trip}}(\beta)=e^{-\beta U}Z_{\mathrm{ph}}(\beta)$
with $Z_{\mathrm{ph}}(\beta)=\Tr e^{-\beta\Omega n_{\mathrm{rel}}}$. Shifting the
energy by $+U$, the remaining three states carry the quadratic form
$H_u=H_0+uW$, $u\in[1,\eta]$, where
\begin{equation}\label{eq:H0W}
H_0=U\mathbf 1_3+\Omega n_{\mathrm{rel}}-\sqrt{U\Omega}\,(P_{20}-P_{02})(a_{\mathrm{rel}}^\dagger+a_{\mathrm{rel}}),
\qquad
W=-\sqrt2\,t\bigl(|\mathrm{sing}\rangle(\langle2,0|+\langle0,2|)+\mathrm{h.c.}\bigr)\otimes\mathbf 1_{\mathrm{rel}} ,
\end{equation}
so that $\cosh(\delta/2)=\eta$ corresponds to $u=\eta$ and no twist to $u=1$.
In the Lang--Firsov frame of the relative mode,
$U_{\mathrm{rel}}=\exp\bigl(-\sqrt X(P_{20}-P_{02})(a_{\mathrm{rel}}^\dagger-a_{\mathrm{rel}})\bigr)$,
one finds $U_{\mathrm{rel}}H_0U_{\mathrm{rel}}^\dagger=UP_{\mathrm{sing}}+\Omega n_{\mathrm{rel}}\ge0$,
so $H_0\ge0$ and $\|W\|_{\mathrm{op}}=2t$; hence
$Z_3(0;\beta)=\Tr_3e^{-\beta H_0}=(2+e^{-\beta U})Z_{\mathrm{ph}}(\beta)>2Z_{\mathrm{ph}}(\beta)$
and, by Weyl's inequality $H_u\le H_0+2ut\mathbf 1_3$,
$Z_3(u;\beta)\ge e^{-2\beta ut}Z_3(0;\beta)>2e^{-2\beta ut}Z_{\mathrm{ph}}(\beta)$.
The triplet weight in $Z_{\mathrm{rel}}(u;\beta)=Z_{\mathrm{trip}}(\beta)+Z_3(u;\beta)$
therefore obeys
\begin{equation}\label{eq:ptrip}
p_{\mathrm{trip}}(u)=\frac{Z_{\mathrm{trip}}(\beta)}{Z_{\mathrm{rel}}(u;\beta)}
\le\frac12e^{-\beta U+2\beta ut}\le\frac12e^{-\frac34\beta U}\le\frac12e^{-3}<0.025
\qquad(\beta U\ge4,\ u\le2,\ t/U\le1/16).
\end{equation}
Let $\Phi_3(u)=\ln Z_3(u;\beta)$. Because $P_{\mathrm{sing}}WP_{\mathrm{sing}}=0$,
$(1-P_{\mathrm{sing}})W(1-P_{\mathrm{sing}})=0$ and $[H_0,P_{\mathrm{sing}}]=0$, one
has $\Tr_3(We^{-\beta H_0})=0$; substituting Duhamel's identity
$e^{-\beta H_u}=e^{-\beta H_0}-u\int_0^{\beta}e^{-(\beta-s)H_u}We^{-sH_0}ds$ into
$\Phi_3'(u)=-\beta\Tr_3(We^{-\beta H_u})/Z_3(u;\beta)$ gives, for every
$u\in[1,\eta]$,
\begin{equation}\label{eq:duhamel}
\Phi_3'(u)=u\beta\int_0^{\beta}\frac{\Tr_3\bigl(e^{-(\beta-s)H_u}We^{-sH_0}W\bigr)}{Z_3(u;\beta)}ds .
\end{equation}
Each integrand is of the form $\Tr_3(XX^\dagger)\ge0$ with
$X=e^{-(\beta-s)H_u/2}We^{-sH_0/2}$, so all of them are real and non-negative and
the integral may be restricted to $s\in[0,s_0]$ with $s_0=1/(100U)\le\beta/400$.
For such $s$ put $\tilde\beta=\beta-s\ge\frac34\beta\ge3/U$ and
$\rho_{3,\tilde\beta}=e^{-\tilde\beta H_u}/Z_3(u;\tilde\beta)$.

Two positivity statements are needed, and one is where the Fock basis would fail.
First, $H_u\ge-2ut\mathbf 1_3$ because $H_0\ge0$, so
$e^{-\beta H_u}=e^{-sH_u}e^{-\tilde\beta H_u}\le e^{2sut}e^{-\tilde\beta H_u}$ and
$Z_3(u;\tilde\beta)/Z_3(u;\beta)\ge e^{-2sut}\ge e^{-1/400}$. Second,
$\Tr_3(\rho_{3,\tilde\beta}W^{2})\ge2t^{2}$. For the second statement work in the
position representation of the relative mode, in which
$a_{\mathrm{rel}}=(x_{\mathrm{rel}}+\partial_{x_{\mathrm{rel}}})/\sqrt2$ and
\begin{equation}\label{eq:posbasis}
\begin{aligned}
H_0&=U\mathbf 1_3+\frac\Omega2\bigl(-\partial_{x_{\mathrm{rel}}}^{2}+x_{\mathrm{rel}}^{2}-1\bigr)
-\sqrt{2U\Omega}\,(P_{20}-P_{02})x_{\mathrm{rel}},\\
uW&=-\sqrt2\,ut\bigl(|\mathrm{sing}\rangle(\langle2,0|+\langle0,2|)+\mathrm{h.c.}\bigr)\otimes\mathbf 1_{L^{2}(\mathbb R)} .
\end{aligned}
\end{equation}
In this basis the kinetic semigroup has a strictly positive Gaussian kernel on
$\mathbb R\times\mathbb R$, the potential is real and diagonal in both the
electronic index and $x_{\mathrm{rel}}$, and $-uW$ has non-negative off-diagonal
electronic entries, so $e^{-\tau uW}$ is entrywise non-negative on $\mathbb C^3$;
by the Trotter product formula the kernel $K_{\tilde\beta}(x,y)_{ij}$ of
$e^{-\tilde\beta H_u}$ is entrywise non-negative. Since
$W^{2}=2t^{2}(\mathbf 1_3+K_{\mathrm{off}})$ with
$K_{\mathrm{off}}=(|2,0\rangle\langle0,2|+|0,2\rangle\langle2,0|+P_{\mathrm{sing}})\otimes\mathbf 1_{L^{2}(\mathbb R)}$
local in $x_{\mathrm{rel}}$ and entrywise non-negative, one gets
$\Tr_3(\rho_{3,\tilde\beta}K_{\mathrm{off}})\ge0$ and hence
$\Tr_3(\rho_{3,\tilde\beta}W^{2})\ge2t^{2}$. In the Fock basis the statement is
false: the term $+\sqrt{U\Omega}(a^\dagger_{\mathrm{rel}}+a_{\mathrm{rel}})$ on
$P_{02}$ has positive off-diagonal entries there, so the oscillator basis is not
the basis in which the semigroup is positivity preserving, and the position basis
is used for that reason.

Two operator identities now convert these into a lower bound for the integrand of
\eqref{eq:duhamel}. Because $W$ commutes with $U\mathbf 1_3+\Omega n_{\mathrm{rel}}$
and $W(P_{20}-P_{02})W=0$,
\begin{equation}\label{eq:WH0W}
WH_0W=W^{2}\bigl(U\mathbf 1_3+\Omega n_{\mathrm{rel}}\bigr)
\end{equation}
exactly on $\mathcal H_3$. Since
$a_{\mathrm{rel}}=\tilde a_{\mathrm{rel}}+\sqrt X(P_{20}-P_{02})$ and
$(A+B)^\dagger(A+B)\le2A^\dagger A+2B^\dagger B$, one has
$\Omega n_{\mathrm{rel}}\le2\Omega\tilde n_{\mathrm{rel}}+2U(P_{20}+P_{02})\le2H_0+2U\mathbf 1_3$.
By convexity of $\beta'\mapsto\ln Z_3(u;\beta')$ on $[\tilde\beta/2,\tilde\beta]$,
\begin{equation}\label{eq:meanH}
\langle H_u\rangle_{\rho_{3,\tilde\beta}}\le\frac2{\tilde\beta}
\ln\frac{Z_3(u;\tilde\beta/2)}{Z_3(u;\tilde\beta)}
\le\frac2{\tilde\beta}\ln\frac{3e^{\tilde\beta ut}Z_{\mathrm{ph}}(\tilde\beta/2)}{2e^{-2\tilde\beta ut}Z_{\mathrm{ph}}(\tilde\beta)}
\le\frac{2\ln3}{\tilde\beta}+6ut ,
\end{equation}
where $Z_3(0;\tilde\beta/2)\le3Z_{\mathrm{ph}}(\tilde\beta/2)$ and
$Z_{\mathrm{ph}}(\tilde\beta/2)/Z_{\mathrm{ph}}(\tilde\beta)=1+e^{-\tilde\beta\Omega/2}<2$,
so $\frac32Z_{\mathrm{ph}}(\tilde\beta/2)/Z_{\mathrm{ph}}(\tilde\beta)<3$. With
$\tilde\beta\ge\frac34\beta\ge3/U$ and $ut\le2U/16$ this gives
$\langle H_0\rangle_{\rho_{3,\tilde\beta}}\le\langle H_u\rangle+2ut\le\frac{2\ln3}{3}U+U<\frac74U$.
Because $0\le W^{2}\le4t^{2}\mathbf 1_3$ commutes with $\Omega n_{\mathrm{rel}}\ge0$,
\begin{equation}\label{eq:WH0W2}
\frac{\Tr_3(\rho_{3,\tilde\beta}WH_0W)}{\Tr_3(\rho_{3,\tilde\beta}W^{2})}
=U+\frac{\Tr_3(\rho_{3,\tilde\beta}W^{2}\Omega n_{\mathrm{rel}})}{\Tr_3(\rho_{3,\tilde\beta}W^{2})}
\le5U+4\langle H_0\rangle_{\rho_{3,\tilde\beta}}\le12U .
\end{equation}
Finally, Jensen's inequality for the convex function $E\mapsto e^{-sE}$ with
respect to the spectral measure of $H_0\ge0$ weighted by
$\rho_{3,\tilde\beta}^{1/2}W(\cdot)W\rho_{3,\tilde\beta}^{1/2}/\Tr_3(\rho_{3,\tilde\beta}W^{2})$
gives, for all $s\ge0$,
\begin{equation}\label{eq:jensen}
\Tr_3\bigl(\rho_{3,\tilde\beta}We^{-sH_0}W\bigr)\ge\Tr_3\bigl(\rho_{3,\tilde\beta}W^{2}\bigr)
\exp\Bigl(-s\frac{\Tr_3(\rho_{3,\tilde\beta}WH_0W)}{\Tr_3(\rho_{3,\tilde\beta}W^{2})}\Bigr)\ge2t^{2}e^{-12sU} .
\end{equation}
Using also $Z_3(u;\tilde\beta)/Z_3(u;\beta)\ge e^{-1/400}$ in
\eqref{eq:duhamel} and integrating over $s\in[0,s_0]$,
\begin{equation}\label{eq:phiprime}
\Phi_3'(u)\ge u\beta\int_0^{s_0}e^{-1/400}\,2t^{2}e^{-12sU}\,ds\ge0.0176\,\beta\,\frac{t^{2}}{U}\,u ,
\end{equation}
because $e^{-1/400-12/100}=e^{-49/400}>0.88$ and $1/100$ enters the length of the
interval. Multiplying by $1-p_{\mathrm{trip}}(u)\ge0.975$ from \eqref{eq:ptrip}
gives $0.975\times0.0176=0.01716>0.01=\frac{2}{200}$, so
$\frac{d}{du}\ln Z_{\mathrm{rel}}(u;\beta)\ge\frac{1}{200}\beta\frac{t^{2}}{U}(2u)$,
and integrating in $u$ from $1$ to $\eta$ yields
\begin{equation}\label{eq:finalgain}
-\bigl(F_{\mathrm{sym}}(\delta)-F_{\mathrm{sym}}(0)\bigr)
\ge\frac{1}{200}\frac{t^{2}}{U}(\eta^{2}-1)=\frac{1}{200}\frac{t^{2}}{U}\sinh^{2}(\delta/2)\ge\frac{1}{400}\frac{t^{2}}{U}\sinh^{2}(\delta/2).
\end{equation}
The derivation produces $1/200$; the last step weakens it to $1/400$ so that the
statement and the certificate carry one and the same constant.
\end{proof}

\subsection{Why this bounds a microscopic extension of the method}

The free-energy gain \eqref{eq:gain} is what a microscopic rotation has to beat.
In the complex-rotation bound the pair field at distance $d$ gains the factor
$(d+1)^{-\kappa}$, and the gain costs the factor $e^{\beta\|C\|_{\mathrm{op}}}$,
where $C$ collects the quadratic part of the twist. Proposition~\ref{prop:floor}
identifies the quadratic part of the twist on a bond, at all phonon frequencies,
and bounds it below by $\frac{1}{400}(t^{2}/U)\sinh^{2}(\delta_e/2)$ with the
constant independent of $U/\Omega$. A rotation which assigns a difference
$\delta_e$ to a bond therefore pays at least
$\beta\frac{1}{400}(t^{2}/U)\sinh^{2}(\delta_e/2)$ on that bond, and the price is
paid again on every further bond across which the profile varies, whereas the pair
field at distance $d$ gains only $\kappa\ln(d+1)$ in total. Since
$\sinh^{2}(\delta_e/2)$ is quadratic in the difference and the difference is
capped at $4\ln2$ by the convergence of the series, the accumulated price grows
with the number of bonds crossed and beats $\kappa\ln(d+1)$ at every distance once
$\beta t^{2}/U$ exceeds a constant multiple of $\kappa$; that scale is the ceiling
of the released $2^{27}t^{2}/U$ bound and not the narrowed ceiling of
Theorem~\ref{thm:main}. What the cold projection does is remove the excursions
$|2,0\rangle\leftrightarrow|\mathrm{sing}\rangle$ from the intermediate states, by
declaring the projector $1-n_{z,\mathrm{odd}}$ to be part of the state; what it
costs is Proposition~\ref{prop:thermal}(2) below.

The honest scope of this section is bond-local. Proposition~\ref{prop:floor} is
proved for every $\Omega>0$, $\beta U\ge4$ and $t/U\le1/16$ on two sites with
$N=2$ and $S_z=0$, and it is that statement which supplies the mechanism. A
lattice-wide lower bound on the twisted free-energy gain which holds for every
rotation profile and every filling is not proved here, so the conclusion drawn
above is a consequence of the mechanism and not a second theorem. What is a
theorem is that no choice of $\Omega$ removes the two-site floor.

\subsection{Numerical cross-check}\label{sec:ed}

The two-site reduction of the proof can be checked independently. Diagonalising
the relative-mode Hamiltonian in the Lang--Firsov basis at a truncated phonon
occupation $N_{\mathrm{ph}}$ and repeating at $N_{\mathrm{ph}}'$,
one obtains the ratio
$R=-(F_{\mathrm{sym}}(\delta)-F_{\mathrm{sym}}(0))\big/\bigl[\frac{t^{2}}{U}\sinh^{2}(\delta/2)\bigr]$,
which is proved to be at least $1/400$, and the positivity statistic
$\Tr_3(\rho_{3,\beta}W^{2})/(2t^{2})$, which is proved to be at least one. The
five rows of Table~\ref{tab:ed} are independent parameter points, deliberately
chosen across both regimes $X<1$ and $X>1$ and across the whole allowed range of
$t/U$ and $\beta U$. The last two columns are the value of $R$ at the two
truncations and their difference, which measures the truncation error, and the
positivity statistic at the finer truncation. Every row satisfies both proved
inequalities with a wide margin, and the truncation error is at most
$2.26\times10^{-6}$. The rows were produced by an independent implementation of
this diagonalisation at the parameters listed; the numerical mode of the
supplement covers the same two inequalities at its own parameter points and
reproduces the crossover reported in Section~\ref{sec:thermal} at the fixed point
used there.

\begin{table}[t]
\centering
\small
\begin{tabular}{@{}rrrrrrrrr@{}}
\hline
$U$ & $\Omega/U$ & $X$ & $\beta U$ & $t/U$ & $\eta$ & $R$ & $|R(N_2)-R(N_1)|$ & $\Tr_3(\rho W^{2})/(2t^{2})$ \\
\hline
$20.93$ & $4.2830$ & $0.2335$ & $20.33$ & $0.0603$ & $1.565$ & $1.875140$ & $4.40\times10^{-14}$ & $1.072769$ \\
$6.58$  & $0.5742$ & $1.7415$ & $16.19$ & $0.0210$ & $1.534$ & $1.156151$ & $5.47\times10^{-12}$ & $1.000737$ \\
$17.53$ & $0.0829$ & $12.0627$ & $20.99$ & $0.0303$ & $1.195$ & $1.031566$ & $2.26\times10^{-6}$ & $1.001012$ \\
$7.51$  & $0.3622$ & $2.7609$ & $9.45$  & $0.0595$ & $1.060$ & $1.110194$ & $9.55\times10^{-12}$ & $1.005070$ \\
$14.52$ & $3.1461$ & $0.3179$ & $16.18$ & $0.0434$ & $1.404$ & $1.637838$ & $1.86\times10^{-12}$ & $1.023963$ \\
\hline
\end{tabular}
\caption{Two-site cross-check of Proposition~\ref{prop:floor}(2). $R$ is the ratio
of the converged free-energy gain to $\frac{t^{2}}{U}\sinh^{2}(\delta/2)$, proved
to be at least $1/400=0.0025$; the last column is the positivity statistic proved
to be at least one. Truncations are $30\to50$ phonons for the first, second and
fifth row, $61\to94$ for the third and $33\to52$ for the fourth.}
\label{tab:ed}
\end{table}

The same diagonalisation at $t=0.5$, $U=20$, $\Omega=1$ reproduces the thermal
crossover of Proposition~\ref{prop:thermal}(1): the Boltzmann-averaged parity
splitting of the two-site spectrum rises from $1.40\times10^{-10}$ at
$\tau=0.5\Omega$ to $2.02\times10^{-2}$ at $\tau=8\Omega$, with the exponential
suppression $e^{-2U/\Omega}$ at low temperature crossing over to the thermal
factor $e^{-U/\tau}$.

\section{Obstruction II: thermal pair transfer and the Gibbs weight of the cold manifold}\label{sec:thermal}

Theorem~\ref{thm:main} is a statement about the Gibbs state of
$H_{\mathrm{cold}}$, and its exponent came from a projection. Two questions arise
about that projection. The first is whether the pair correlations it discards are
the ones that decay slowly; the answer is that in the strong-coupling regime the
thermal pair transfer is itself small at low temperature but crosses over to a
thermal factor at temperatures above the phonon frequency. The second is whether
the projection can be removed by restricting the Gibbs trace to the cold manifold;
the answer is no, since the weight of the cold manifold is exponentially small in
the volume at every fixed temperature. Both statements are elementary, and both
are needed to state the open problem honestly.

\begin{proposition}[Thermal pair transfer and the cold weight]\label{prop:thermal}
\begin{enumerate}
\item\label{it:hky} In the effective thermal pair hopping of
\cite[Appendix D]{hky}, $t_b(\tau)\sim\frac{t^{2}}{U}\exp\bigl(-\frac{2U}{\Omega}\tanh(\beta\Omega/2)\bigr)$,
the exponent crosses over from $2U/\Omega$ for $\tau\ll\Omega$ to $U/\tau$ for
$\tau\gg\Omega$, so that $\Omega\ll\tau\ll U$ gives
$t_b(\tau)\sim\frac{t^{2}}{U}e^{-U/\tau}$. Moreover, for all $t,U>0$,
\begin{equation}\label{eq:peak}
\max_{\tau>0}\ \beta\,\frac{t^{2}}{U}\,e^{-U/\tau}=\frac{t^{2}}{e\,U^{2}},
\qquad\text{attained at }\tau=U .
\end{equation}
\item\label{it:weight} For the atomic oscillator Hamiltonian
$H_{\mathrm{at,LF}}=\frac U2\Nodd+\Omega\sum_{z}n^{\mathrm{ph}}_z-\mu N$, at every
$\mu\in\mathbb R$, every $\beta>0$ and every $L\ge4$,
\begin{equation}\label{eq:weight}
\begin{aligned}
\frac{\Tr\bigl(P_{\mathrm{cold}}e^{-\beta H_{\mathrm{at,LF}}}\bigr)}
{\Tr\bigl(e^{-\beta H_{\mathrm{at,LF}}}\bigr)}
&=\Bigl[\frac{(1+e^{2\beta\mu})\bigl(1-e^{-\beta\Omega(M+1)}\bigr)}
{1+2e^{\beta\mu}e^{-\beta U/2}+e^{2\beta\mu}}\Bigr]^{|\Lambda|}\\
&\le\bigl(1-e^{-\beta\Omega(M+1)}\bigr)^{|\Lambda|}
\le\exp\bigl(-|\Lambda|\,e^{-\beta\Omega(M+1)}\bigr),
\end{aligned}
\end{equation}
and at $\mu=0$ the middle ratio reads
$\bigl[\frac{2(1-e^{-\beta\Omega(M+1)})}{2+2e^{-\beta U/2}}\bigr]^{|\Lambda|}$.
\end{enumerate}
\end{proposition}

\begin{proof}
For part~\ref{it:hky}, the exponent is $2U\Omega^{-1}\tanh(\beta\Omega/2)$, which
equals $2U/\Omega$ up to exponentially small corrections $O(e^{-\beta\Omega})$
when $\beta\Omega\gg1$ and equals $\beta U=U/\tau$ up to relative corrections of
order $(\beta\Omega)^{2}$ when $\beta\Omega\ll1$. For the peak,
substitute $z=U/\tau>0$, so that $\beta e^{-U/\tau}=\tau^{-1}e^{-U/\tau}=(z/U)e^{-z}$,
which is maximal at $z=1$ with value $1/(eU)$; multiplying by $t^{2}/U$ gives
\eqref{eq:peak}. The maximum identity is certified in exact arithmetic.

For part~\ref{it:weight}, $H_{\mathrm{at,LF}}$ is a sum of single-site operators,
so both traces factorise over sites. At one site the electronic configurations are
$n=0$, $n=2$ (both with $\Nodd=0$) and the two singly occupied states, and the
trace over electronic states of
$e^{-\beta((U/2)\Nodd-\mu n)}$ times the phonon trace gives
\begin{equation}\label{eq:site}
\Tr_{\mathrm{site}}\bigl(e^{-\beta H_{\mathrm{at,LF}}}\bigr)
=\bigl(1+e^{2\beta\mu}+2e^{\beta\mu}e^{-\beta U/2}\bigr)\sum_{n\ge0}e^{-\beta\Omega n}
=\frac{1+e^{2\beta\mu}+2e^{\beta\mu}e^{-\beta U/2}}{1-e^{-\beta\Omega}}
\end{equation}
while the projected trace keeps only the occupations $n\le M$ at sites with
$\Nodd=0$, which gives
\begin{equation}\label{eq:siteproj}
\Tr_{\mathrm{site}}\bigl(P_{\mathrm{cold}}e^{-\beta H_{\mathrm{at,LF}}}\bigr)
=\bigl(1+e^{2\beta\mu}\bigr)\sum_{n=0}^{M}e^{-\beta\Omega n}
=\frac{(1+e^{2\beta\mu})\bigl(1-e^{-\beta\Omega(M+1)}\bigr)}{1-e^{-\beta\Omega}} .
\end{equation}
The ratio of \eqref{eq:siteproj} to \eqref{eq:site} is the bracket in
\eqref{eq:weight}, and it is at most $1-e^{-\beta\Omega(M+1)}$ because
$1+e^{2\beta\mu}\le1+e^{2\beta\mu}+2e^{\beta\mu}e^{-\beta U/2}$. Raising to the
power $|\Lambda|$ and using $\ln(1-x)\le-x$ for $x\in(0,1)$ gives the two bounds.
\end{proof}

The content of part~\ref{it:weight} is geometric in the volume at fixed
temperature. Since $M=\lfloor X/64\rfloor$ and $X=U/\Omega$, one has
$\beta\Omega(M+1)\le\beta U/64+\beta\Omega$, so
\begin{equation}\label{eq:weightbound}
\frac{\Tr(P_{\mathrm{cold}}e^{-\beta H_{\mathrm{at,LF}}})}{\Tr(e^{-\beta H_{\mathrm{at,LF}}})}
\le\exp\bigl(-|\Lambda|\,e^{-\beta U/64-\beta\Omega}\bigr),
\end{equation}
which is exponentially small in $|\Lambda|$ for every fixed $\tau>0$. Restricting
the Gibbs trace of the full model to the cold manifold is therefore not an
approximation that improves with volume; it discards almost all of the state. The
projection is a modelling choice, made explicit in Section~2.3, and the pair
correlations of Theorem~\ref{thm:main} are those of the projected model. Whether
the same decay with the same exponential narrowing holds for the Gibbs state of
the full Holstein Hamiltonian at phonon occupations above $M$ is the open
question of this note.

\section{Limitations}\label{sec:limits}

The following restrictions should be read together with the statements.

\emph{The model is an effective model.} Theorem~\ref{thm:main} concerns
$H_{\mathrm{cold}}$ of \eqref{eq:Hcold}, that is, the Holstein Hamiltonian after a
Lang--Firsov transformation, a first-order Schrieffer--Wolff truncation and a
projection onto $\Nodd=0$ with phonon occupations at most $M$. The truncation
discards the second and higher orders of the generator, which carry their own
smallness in $x=t/U\le2^{-14}$ but are not controlled here beyond the first order;
and the projection is a restriction of the state, with the weight computed in
Proposition~\ref{prop:thermal}(2). No statement is made about the Gibbs state of
\eqref{eq:H} itself.

\emph{The full-Gibbs control for $n^{\mathrm{ph}}>M$ is open.} The exponential
narrowing $e^{-7U/(8\Omega)}$ is produced by the Bargmann--Fock bound
\eqref{eq:bf} applied between phonon occupations $M+1$ and $M$. Admitting
occupations above $M$ in the Gibbs state reintroduces the neutral excursions whose
free-energy gain was bounded below in Proposition~\ref{prop:floor}, and the two
mechanisms have not been reconciled in a single estimate. A proof that the full
Gibbs pair correlations obey an exponentially narrowed ceiling in $U/\Omega$ (or
a sub-$t^{2}/U$ ceiling such as $C\Omega$), or that they do not, is not given.

\emph{The bounds are on a clustering temperature, not on a critical temperature.}
$T_{\mathrm{pair}}$ and $T^{\mathrm{cold}}_{\mathrm{pair}}$ are thresholds for the
decay and summability of one two-point function, as in
Definition~\ref{def:tpair}. No phase transition is proved or excluded, at any
temperature, in either the full or the projected model, and nothing here says that
the model is a superconductor or possesses any form of long-range order.

\emph{The prefactor is not optimised.} The constant $2^{26}$ in \eqref{eq:ceiling}
is the power of two that dominates the certified value
$64K_4(x)/x^{2}<5.562\times10^{7}$, whose logarithm to base two is $25.73$. The
series \eqref{eq:K4def} was not minimised over the choices made in the counting,
and the exponent $\xi_{\mathrm{cold}}$ is the value produced by the rotation
profile \eqref{eq:theta} with the cap $\kappa\le4$; neither is claimed to be
optimal. The comparison with the released full-Gibbs ceilings $2^{27}t^{2}/U$
\cite{rens2026a} and $6000\,t^{2}/U$ \cite{rens2026strong} is a comparison of
numbers obtained by different constructions and with different scopes: both
released bounds hold for the full Gibbs state with no $t/U$ cutoff, uniformly in
$\Omega>0$, and yield exponential decay, whereas Theorem~\ref{thm:main} holds
only for the cold-projected model at $t/U\le2^{-14}$ and $X=U/\Omega\ge1$ with
algebraic decay. Numerically, $2^{26}e^{-7X/8}$ lies below $2^{27}$ for all
$X\ge1$ and below $6000$ only for $X>\frac{8}{7}\ln(2^{26}/6000)\approx10.65$,
and neither comparison is an improvement of the released full-Gibbs theorems.

\emph{The obstruction is bond-local.} Proposition~\ref{prop:floor} is a two-site
statement, with $N=2$, $S_z=0$ and $\mu=0$. It isolates the mechanism that the
cold projection suppresses, and it is proved for every $\Omega$, but it is not a
lattice theorem about all rotation profiles. In particular, the statement that a
direct microscopic rotation cannot reach a ceiling below a constant multiple of
$t^{2}/U$ is an inference from that mechanism, and the proved content is the
two-site floor and its numerical confirmation in Table~\ref{tab:ed}.

\emph{What the certificate program does and does not do.} It recomputes the
finite-dimensional operator norms on the sixteen electronic states of a bond, the
coherent-state bounds in the form of outward-rounded intervals, the summation of
the two series in \eqref{eq:K4def} plus a geometric tail, the shell and geometric
sums of Lemma~\ref{lem:shell} and Corollary~\ref{cor:orig}, and the two-site
spectra. It does not formalise the operator inequalities of
Sections~\ref{sec:proof} to \ref{sec:floor}, each of which rests on the written
argument.

\section{Supporting material}\label{sec:supp}

The supplement contains the certificate program \texttt{replay.py} and a list of
file digests. The program was run on the archived copy that is distributed in the
supplement, and the digests in the supplement are the digests of that copy. The program has three modes. Its default run checks five groups of
exact statements in rational arithmetic or in outward-rounded interval arithmetic:
the sixteen-state operator identities on a bond, the Bargmann--Fock and Poisson
bounds, the torus shell telescoping together with the two series of
\eqref{eq:K4def} and the prefactor comparison $64K_4(x)/x^{2}<2^{26}$, the rotated
cloud summability $(1-r)^{-2}<2$ of Corollary~\ref{cor:orig}, and the two-site
bounds of Proposition~\ref{prop:floor}. Its mutant mode corrupts the claimed
inequalities one at a time, sixteen variants in all, and requires each one to be
rejected. Its numerical mode reproduces the thermal crossover reported in
Section~\ref{sec:thermal} at $t=0.5$, $U=20$, $\Omega=1$, and checks the two
inequalities of Proposition~\ref{prop:floor} at six parameter points of its own,
which are not the rows of Table~\ref{tab:ed}. Table~\ref{tab:constants} lists the
constant chain of Theorem~\ref{thm:main} with the place where each constant enters.

\begin{table}[t]
\centering
\small
\begin{tabular}{@{}L{0.34\textwidth}L{0.30\textwidth}L{0.24\textwidth}@{}}
\hline
Quantity & Value & Where \\
\hline
$\cosh(2\ln2)$ & $17/8$ & \eqref{eq:s1norm} \\
$\|S_{1,e}\|_{\mathrm{op}},\ \|\ad_{S_{1,e}}\|_{\mathrm{op}}$ &
$\frac{17\sqrt2}{4}x<7x$, $<14x$ & \eqref{eq:s1norm}, \eqref{eq:K4def} \\
$\|V_{e,\mathrm{mix}}\|_{\mathrm{op}}+\|V_{e,\mathrm{diag}}\|_{\mathrm{op}}$ &
$<7t+5t=12Ux$ & \eqref{eq:K4} \\
coherent-state factor at two sites & $(\sqrt2)^{2}e^{-7X/8}=2e^{-7X/8}$ & \eqref{eq:bf} \\
chain ratio & $r=4\cdot14x=56x\le7/2048$ & \eqref{eq:famIIterm} \\
series & $K_4(x)\le2^{20}x^{2}$ & \eqref{eq:K4def} \\
prefactor & $64K_4(x)/x^{2}<5.562\times10^{7}<2^{26}$ & \eqref{eq:ceiling} \\
exponent & $\xi_{\mathrm{cold}}=4-16A$ or $1/(4A)$ & \eqref{eq:decay} \\
cold-correlation prefactor & $4=2^{2}$, from $(1-r)^{-2}<2$ twice & Corollary~\ref{cor:orig} \\
\hline
\end{tabular}
\caption{The constants of Theorem~\ref{thm:main} and Corollary~\ref{cor:orig},
with the place of entry. The values in the second column are certified in exact
arithmetic at $x=2^{-14}$ and $X=1$.}
\label{tab:constants}
\end{table}

\section*{Records}

The strong-coupling companion preprint on the ceiling $6000\,t^{2}/U$ is
\cite{rens2026strong}, \record{\STRONGDOI}{whose record is reserved for release};
its replay supplement is archived at \DOI{10.5281/zenodo.24414618}. The earlier release with the uniform
ceiling $2^{27}t^{2}/U$ is \cite{rens2026a}, and the companion analysis of
hypercubic lattices in three and more dimensions is \cite{rens2026b}. The present
note is archived at \record{\SELFDOI}{the record reserved for it}, and its
supplement is archived together with it.

\section*{AI assistance}

The analytical results, the design of the estimates, and the final text are the
author's responsibility. Independent proof reviews of the argument and of the
certificate program were carried out by separate AI systems, whose comments led
to several corrections recorded in the present version. The manuscript, its LaTeX
source, and the certificate program were drafted and edited with AI assistance
under the author's direction and review.
\addcontentsline{toc}{section}{AI assistance}

This preprint has not been peer reviewed.

\bibliographystyle{plain}
\bibliography{refs}

\end{document}
